\documentclass[10pt,letterpaper]{article}
\usepackage[T1]{fontenc}
\usepackage{lmodern,amsmath,amssymb,amsthm,mathtools}
\usepackage[margin=0.85in]{geometry}
\usepackage{microtype,booktabs,longtable,array,enumitem,xcolor,fancyhdr}
\usepackage{hyperref,xurl}
\hypersetup{colorlinks=true,linkcolor=blue!45!black,urlcolor=blue!45!black,
 pdftitle={Cofinal Pell Centers and Oriented Power Certificates},
 pdfauthor={AI-assisted research prepared for Vladimir Reshetnikov}}
\setlength{\headheight}{14pt}
\pagestyle{fancy}\fancyhf{}
\fancyhead[L]{\small Explicit Diophantine forms of Con(ZFC)}
\fancyhead[R]{\small 25 October 2026}\fancyfoot[C]{\thepage}
\setlength{\parskip}{4pt}\setlength{\parindent}{0pt}
\setlist{nosep,leftmargin=*}
\allowdisplaybreaks
\newtheorem{theorem}{Theorem}[section]
\newtheorem{lemma}[theorem]{Lemma}
\newtheorem{proposition}[theorem]{Proposition}
\newtheorem{corollary}[theorem]{Corollary}
\theoremstyle{definition}\newtheorem{definition}[theorem]{Definition}
\theoremstyle{remark}\newtheorem{remark}[theorem]{Remark}
\newcommand{\Z}{\mathbb Z}\newcommand{\N}{\mathbb N}\newcommand{\Q}{\mathbb Q}
\newcommand{\Con}{\operatorname{Con}}\newcommand{\ZFC}{\mathsf{ZFC}}
\newcommand{\Prf}{\operatorname{Prf}}
\newcommand{\pin}{\texttt{02a78a08603cbae248c19c23a75907c6c6bd6443}}
\title{\textbf{Cofinal Pell Centers and Oriented Power Certificates}\\[4pt]
\large Four explicit Diophantine equations for syntactic ZFC consistency}
\author{AI-assisted research prepared for Vladimir Reshetnikov\\
\small An additive contribution to ProveIt}
\date{25 October 2026}
\begin{document}
\maketitle
\begin{abstract}
We give two arithmetic reductions and four fully specified integer
polynomials whose absence of integer zeros is provably equivalent in ZFC
to its syntactic consistency, by composition with the identified ProveIt
C222 refutation calibration. First, an even Pell frame permits the free
center $A=L^2+2u^2$, replacing $L^2+u^2+v^2$. The new ingredient is an
order-eight identity in a finite quadratic algebra. Canonical parity in
the existing source makes the required frame change only a constant
increment. The resulting quartics have 229 and 228 integer variables,
respectively 4,477 and 12,001 collected monomials, and coefficient height
7,968. Second, oriented reciprocal pooling combines power certificates
at different, initially signed bases. Its conclusion applies to nonzero
bases while zero bases remain inactive; a parity argument then proves
that every actual base is nonzero without circular reasoning. This gives
206 variables at degree 24 and 207 at degree 16. All variables are
unrestricted integers. Literal residuals, exact coefficients,
reconstruction programs, and independent finite checks accompany the
proofs. These are measured improvements over named parents and new
tradeoffs; a global simplest equation, worldwide priority, and a new
end-to-end proof-assistant derivation are not claimed.
\end{abstract}
\setcounter{tocdepth}{1}
\tableofcontents
\newpage

\section{Result, specification, and comparison}

\subsection{The precise logical target}

Fix the effective first-order ZFC presentation called
\texttt{ZFCax\_s} in ProveIt and its finite derivation relation
\texttt{BProv}. Let $\Prf_{\ZFC}(p,\bot)$ be its arithmetized proof
predicate. Throughout,
\[
 \Con(\ZFC)\;\equiv\;\neg\exists p\in\N\;\Prf_{\ZFC}(p,\bot).
\]
Effective translations between ordinary presentations give the usual
syntactic consistency statement. No assumption about externally standard
integers, transitive set models, or the truth of all ZFC theorems is
substituted for this definition.

For each of the four polynomials below the result is
\begin{equation}\label{eq:target}
 \ZFC\vdash
 \left(\exists\boldsymbol{x}\in\Z^n\;P(\boldsymbol{x})=0\right)
 \ \longleftrightarrow\
 \exists p\in\N\;\Prf_{\ZFC}(p,\bot).
\end{equation}
Negating both sides gives the requested no-solution equivalence. The
claim is an ordinary mathematical proof formalizable in ZFC, using the
explicitly credited inherited calibration; the archive does not contain
a newly compiled Lean/Rocq proof of the whole statement or a literal
object-language ZFC derivation. It proves neither side of the consistency
question.

\subsection{The delivered polynomials}

Every polynomial is specified by a finite list of integer-coefficient
residuals and the instruction to square and sum that list. All named
expressions are polynomial abbreviations, not additional quantified
variables. The JSON convention is
\[
 [c,[i_1,\ldots,i_d]]\quad\longleftrightarrow\quad
 c\,x_{i_1}\cdots x_{i_d};\qquad[c,[]]\longleftrightarrow c.
\]
Repeated indices are fixed nonnegative exponents. No exponentiation
primitive, factorial, binomial primitive, free coefficient parameter,
congruence side condition, or implicit sign restriction survives in a
final equation.

\begin{center}\small
\begin{tabular}{@{}lrrrrr@{}}
\toprule
Polynomial & Variables & Degree & Residuals & Monomials & Height\\
\midrule
$P_{229}$, balanced &229&4&176&4,477&7,968\\
$P_{228}$, compact &228&4&175&12,001&7,968\\
$P_{207}$, split pooling &207&16&146&94,294&13,607,065,661\\
$P_{206}$, full pooling &206&24&145&not fully expanded&not claimed\\
\bottomrule
\end{tabular}
\end{center}
``Height'' means the maximum absolute coefficient of the fully collected
single polynomial. For $P_{206}$ the complete residual list is supplied:
144 quadratic residuals and one degree-12 residual. The latter has
25,144 monomials. Its square need not be expanded to specify the
polynomial completely. We deliberately do not report a collected
coefficient height for that unperformed expansion.

$P_{207}$ has the same 144 quadratic residuals, a degree-8 residual and a
degree-5 residual. Its full expansion is included. The maximum absolute
coefficient within the residual lists of both pooled equations is
95,232; this is a different statistic from final polynomial height.

\subsection{What the improvements measure}

The comparison snapshot is \pin. It was the observed \texttt{main} head
when the source review began. We inspected the three requested repository
locations, the active research indexes, and the later delivered packages;
the collection's introductory frontier table lags several of its newest
files. In particular, package 173 already supplies a 215-variable
quartic, and 213/214-variable quartics with much larger coefficients.
We do not compare against obsolete 225-variable records as though those
were the whole current frontier.

\begin{center}\small
\begin{tabular}{@{}lrrrr@{}}
\toprule
Matched quartic & Variables & Monomials & Height & Residuals\\
\midrule
Package 166 balanced parent &230&4,482&7,968&176\\
New balanced center &229&4,477&7,968&176\\
Package 166 compact parent &229&12,006&7,968&175\\
New compact center &228&12,001&7,968&175\\
\bottomrule
\end{tabular}
\end{center}
Each center reduction saves exactly one integer coordinate and five
collected monomials, with degree, height, and residual count unchanged.
The maximum residual coefficient also remains 64. These are strict
improvements on the two actual parent equations, not a new record for
quartic variable count.

Pooling starts from package 173's 221-coordinate, 160-row source chart,
the source of its 215-variable balanced quartic. Replacing ten private
certificate coordinates by one or two saves nine or eight coordinates;
the same six certified projections then give 206 or 207. The degree
increases, so these points do not dominate the existing quartics. The
repository also contains nine-positive and ten-integer constructions,
and a 37-integer quartic with a very different coefficient architecture
\cite{collection,sun}. ``Simplest'' therefore requires choosing a
measure: variables, degree, support, coefficient bits, evaluation cost,
proof complexity, or some specified combination.

The most reusable contribution here is the arithmetic behind the
reductions. The finite quadratic-ring construction is inherited from
package 166, but the order-eight square correction strengthens its
two-square conclusion. Package 168 already pointed toward multi-base
reciprocal pooling; the oriented, zero-base-tolerant theorem below and
its noncircular activation argument complete a concrete version of that
direction. Literature-wide novelty is not established.

\section{One scaled square for a cofinal Pell center}

For $a\ge2$ and $j\ge0$ define the standard Pell polynomials by
\[
 (a+\sqrt{a^2-1})^j=X_j(a)+Y_j(a)\sqrt{a^2-1}.
\]
Equivalently, $X_0=1$, $Y_0=0$, $X_1=a$, $Y_1=1$, and either sequence
satisfies $U_{j+2}=2aU_{j+1}-U_j$. The identities used below follow by
polynomial multiplication. The finite quadratic algebra notation is
only a proof device; it does not enlarge the final polynomial's domain.

\begin{theorem}[One scaled square]\label{thm:center}
Let $W\ge2$ be even. Let $q,M$ be positive odd integers, with every
prime divisor of $M$ dividing $q$. Write $n=q+1$, and put
\[
 a_0\equiv\frac{n^2+1}{2n}
       =1+\frac{q^2}{2(q+1)}\pmod M.
\]
Then there exist an odd positive integer $k$ and an integer $u$ such that
\begin{equation}\label{eq:onecenter}
 L=X_{Wk}(W+1),\qquad A=L^2+2u^2,\qquad A\equiv a_0\pmod M.
\end{equation}
For this fixed $k$, the center $A$ can be made arbitrarily large.
Moreover $Y_{Wk}(W+1)/W$ is an odd positive integer. Thus, for some
integer $z$, the inherited frame equation is preserved:
\begin{equation}\label{eq:oddframe}
 L^2-W^3(W+2)(2z+1)^2=1.
\end{equation}
\end{theorem}

\begin{proof}
If $M=1$, choose $k=1$ and let $u$ grow; the parity assertion follows
from the argument at the end. Otherwise put $D=W(W+2)$ and work in
\[
 S_M=(\Z/M\Z)[s]/(s^2-D).
\]
The element $e=(W+1)+s$ has norm one and inverse $(W+1)-s$. Its
$W$th power has finite order $2^r k$, with $k$ odd. Therefore
$g=e^{Wk}$ has power-of-two order. Write $g=L+Ys$. These are precisely
the integer Pell coordinates in \eqref{eq:onecenter}, reduced modulo $M$.

Fix $p^a\Vert M$, where $p$ is odd. We shall prove that
\[
 c=\frac{a_0-L^2}{2}\pmod {p^a}
\]
is a square. Inverses of all displayed denominators exist because
$p\mid q$ and $p\nmid2(q+1)$.

\emph{Case 1: $L\equiv\epsilon\pmod p$, with $\epsilon=\pm1$.}
If $D\not\equiv0\pmod p$, the norm relation
$L^2-DY^2=1$ gives $Y=0\pmod p$. If $D=0\pmod p$, then $s^2=0$.
Raising $\epsilon+Ys$ to a power of two shows that its nilpotent
coefficient must vanish: that coefficient is an invertible multiple of
$Y$ over the odd-characteristic field. Thus $g=\epsilon$ in $S_p$
in either case.

The kernel of reduction from $S_{p^a}^{\times}$ to $S_p^{\times}$
consists of elements $1+pv$ and is a finite $p$-group. The element
$\epsilon g$ belongs to this kernel and has power-of-two order, so it
is the identity. Hence $L^2=1\pmod{p^a}$ and
\[
 c\equiv\frac{q^2}{4(q+1)}\pmod {p^a}.
\]
The unit $q+1$ is $1$ modulo $p$. The root $1$ of
$b^2-(q+1)$ has derivative $2$ modulo $p$ and lifts successively to a
unit root $b$ modulo $p^a$. Thus $c=(q/(2b))^2$ is a square.

\emph{Case 2: $L\not\equiv\pm1\pmod p$.}
The ramified case $D=0\pmod p$ is impossible by the same nilpotent
argument. Also $Y\ne0\pmod p$. In a separable quadratic algebra over
$\mathbb F_p$, the only norm-one elements with square one are $1$ and
$-1$: this is immediate in a field, and in the split algebra the two
coordinates are mutual inverses. Therefore the image of $g$ has order
at least four. Because $W$ is even and $k$ odd, the order of $e$ modulo
$p$ is divisible by eight.

Choose a power $h=H+Zs$ of $e$ of exact order eight. The unique
nonidentity norm-one involution is $-1$, so $h^4=-1$. Hence
$h^{-2}=-h^2$, and comparison of scalar parts gives
\[
 2H^2=1,\qquad DZ^2=H^2-1=-\tfrac12,
 \qquad -\frac D2=(DZ)^2
 \quad\text{in }\mathbb F_p.
\]
Since $a_0=1\pmod p$,
\[
 c=\frac{1-L^2}{2}=-\frac{DY^2}{2}=(DZY)^2\pmod p.
\]
It is nonzero, so its square root lifts to modulo $p^a$ by the
invertibility of $2u$. This proves the local claim without quadratic
reciprocity or a formula for the order of the norm-one group.

The Chinese remainder theorem combines the roots into $u\pmod M$.
Replacing $u$ by $u+jM$ preserves the congruence and makes
$L^2+2u^2$ arbitrarily large. In all cases $A\ge L^2$.

It remains to prove the frame parity. At $a=W+1$, the recurrence modulo
$W$ gives $Y_j(a)\equiv j\pmod W$, so $W\mid Y_W(a)$. Because $W$
is even, $a$ is odd and every $X_j(a)$ is odd. The odd-index ordinates
are odd; the doubling formula
$Y_{2j}(a)=2X_j(a)Y_j(a)$ therefore gives
$\nu_2(Y_W(a))=\nu_2(W)$. The quotient $Y_W(a)/W$ is odd. Finally,
\[
 Y_{Wk}(a)=Y_W(a)Y_k(X_W(a)),
\]
and the second factor is odd for odd $k$. The positive quotient is
$2z+1$ for an integer $z$. Substitution in the Pell norm proves
\eqref{eq:oddframe}.
\end{proof}

\subsection{Why the hypotheses matter}

The theorem does not say that an arbitrary prescribed center gap is a
square or that every preselected frame works. The odd multiplier $k$
is chosen before the correction $u$. It is an existence choice for
already present variables $L,z$, not a new quantified polynomial
coordinate.

There are small exact obstructions to tempting stronger statements.
For $W=5,q=3,M=3$, we have $D=35=-1\pmod3$ and $e=s$. For every
odd $k$, $e^{5k}$ has scalar coordinate $L=0\pmod3$, whereas
$a_0=1\pmod3$. But $2u^2$ is never $1\pmod3$. Thus the even-frame
hypothesis cannot simply be discarded.

For $W=6,q=3,M=27$, every index $6k$ has $L^2=1\pmod{27}$:
in the quadratic algebra, $(7+s)^6=1+12s$ and $(12s)^2=0$ modulo 27,
so every power of this element has scalar part one.
Here $a_0=19\pmod{27}$. The required correction $18\pmod{27}$
is not a square, but it equals $2\cdot3^2$. A single unscaled square
would fail. These examples are also executable negative controls.

\section{Applying the center theorem to the literal quartics}

\subsection{The inherited modular interface}

Package 166 directly compiles a numerical proof calculus. Its six exact
power requests are
\[
 R=16^t,\quad N=B^x,\quad q=(BH)^x,\quad
 D_t=q^t,\quad D_y=q^{y+1},\quad M_{\rm pow}=q^x.
\]
Eleven other numerator values are supplied by modular Pell evaluations
at base $q+1$. They need a common center satisfying
\[
 2(q+1)A-(q+1)^2-1\equiv0\pmod\Omega,
 \qquad \Omega=NM_{\rm pow}.
\]
In completeness, $B,q,\Omega$ are positive odd integers and every prime
of $\Omega$ divides $q$. Thus the residue is exactly that of
Theorem~\ref{thm:center}. The existing two-square theorem supplies
$A=L^2+u^2+v^2$; our theorem supplies $A=L^2+2u^2$ instead.
The signed Pell recognition and common-root modular numerator arguments
are inherited \cite{source166}. Their source proofs are retained in
the archive with their scope and provenance.

\subsection{Making the frame even by one constant}

Write $\Sigma$ for the sum of the 21 atom squares in the source frame.
The atoms, grouped only for readability, are
\[
 t,x,y+1,R+t,B;\qquad a_0,\ldots,a_8;\qquad
 N,J,q,D_t,D_y,M_{\rm pow},BH.
\]
Here $a_0,\ldots,a_8$ denote actual CRT representatives; the notation
$a_0$ in this paragraph is unrelated to the center residue in
Theorem~\ref{thm:center}.

The inherited completeness construction already chooses $x,J$ even,
$R,y$ even, one odd CRT representative and eight even ones, and the
seven quantities $B,N,q,D_t,D_y,M_{\rm pow},BH$ odd. Its later CRT
adjustments preserve these parities. The two atoms $t,R+t$ have the same
parity, whatever the parity of $t$. Therefore
\[
 \Sigma\equiv
 \underbrace{t^2+(R+t)^2}_{0}
 +\underbrace{x^2+J^2}_{0}
 +\underbrace{(y+1)^2+7}_{0}
 +\underbrace{\sum_{i=0}^8a_i^2}_{1}
 \equiv1\pmod2.
\]
Changing the frame from $W=2+\Sigma$ to $W=3+\Sigma$ thus makes
the canonical completeness frame even. All inherited lower bounds are
strengthened. On an arbitrary polynomial zero $W$ need not be even:
soundness only uses the stronger magnitude bound and $A\ge L^2$.
Evenness is used when constructing witnesses, exactly where the parity
calculation applies.

The delivered parity checker reconstructs the complete square sum from
the literal raw source. In its stored coordinates the frame constant
changes from $-525$ to $-526$. It independently propagates the addition
chain for $J$ modulo two and checks both possible parities of $t$.
Finite checking authenticates these identities; availability of the
canonical semantic choices is the inherited completeness theorem.

\subsection{Soundness, completeness, and exact algebra}

For soundness, the new center row gives $A\ge L^2$ and the new frame
implies every previous frame inequality. All retained arithmetic
recognition lemmas apply. The modular numerators, exact binomial
coefficients, finite rule rows and Horn-flow certificate decode as in
the parent, yielding a C222 refutation.

For completeness, start from a finite C222 refutation. Choose the
inherited semantic data with the stated parities, and compute the new
even $W=3+\Sigma$. Apply Theorem~\ref{thm:center} with modulus
$\Omega=Nq^x$. Choose $A$ large
enough for all inherited cap and orientation constructions, then fill
their Pell coordinates, modular numerators, and remaining arithmetic
gates. This constructs a zero of the modified source. The constrained
eliminations and shears have recorded integral restoration maps on their
zero sets; deleted unused coordinates can be restored arbitrarily.

The reconstruction replays the parent source's 208 monic eliminations,
its integral translations, and the balanced or compact final maps.
Every pivot is checked to be an equation $v=f$ with integral $f$
independent of $v$. Each residual remains quadratic. The now-unused
private $v$ is deleted, and every surviving coordinate occurs.

There is also a much shorter independent description at the final
endpoint. Let $F$ be the parent's final frame residual after the substitution
$v=u$ (in fact it is independent of these center coordinates), and identify its
private center coordinates $v=u$. Every residual is then unchanged
except $F\mapsto F-1$. Consequently the complete single polynomials
satisfy the exact identity
\begin{equation}\label{eq:globalidentity}
 \boxed{P_{\rm new}=P_{\rm parent}[v:=u]-2F+1.}
\end{equation}
This identity checks the emitted coefficients. It does not by itself
prove existential equivalence: an arbitrary old center need not have
$u=v$. The cofinal theorem supplies the missing completeness argument
by choosing the frame and center afresh.

\section{Oriented reciprocal pooling at different bases}

The next reduction is independent of the center construction. It
starts from package 173's unchanged signed Pell source, whose frame
already gives $|A|>W^W$ and bounds the squares of every actual base,
requested output, and nominal index by $W$ \cite{source173}.
The Pell polynomials are evaluated at signed integer $A$ using their
polynomial recurrences. For signed indices use $X_{-k}(A)=X_k(A)$ and
$Y_{-k}(A)=-Y_k(A)$. These are the coefficients of the corresponding
positive or negative powers of the unit $A+s$, $s^2=A^2-1$.

\begin{theorem}[Oriented pooling with inactive zero bases]\label{thm:pool}
Let $W\in\Z$, $W\ge258$, let $A\in\Z$ satisfy $|A|>W^W$, and let at most five
integer bases $n_j$ be given. There are at most 64 requests in total.
Each request in block $j$ consists of integers $(m_i,k_i)$ and canonical
signed Pell coordinates $(X_i,Y_i)$ at parameter $A$ and index $k_i$.
Assume
\[
 n_j^2<W,\qquad m_i^2<W,\qquad k_i^2<W.
\]
Define
\begin{align*}
 T_j&=2An_j-n_j^2-1,&D_i&=X_i-AY_i,\\
 S_j&=\sum_{i\in j}\bigl((m_i-D_i)^2-Y_i^2\bigr),&
 H_j&=\sum_{i\in j}Y_i(m_i-X_i),\\
 E_j&=S_j-2n_jH_j.
\end{align*}
All $T_j$ are nonzero. Then
\begin{equation}\label{eq:orientedintegrality}
 \sum_j\frac{n_jE_j}{T_j}\in\Z
\end{equation}
if and only if every request in every nonzero-base block satisfies
$m_i=n_j^{k_i}$, initially interpreted in $\Q$. For $|n_j|\ge2$ this
also forces $k_i\ge0$. No conclusion is asserted about a zero-base
block's outputs; its term in \eqref{eq:orientedintegrality} is zero.
\end{theorem}

\begin{proof}
For $n_j=0$, $T_j=-1$. For $n_j\ne0$,
\[
 |T_j|\ge2|A|-W>W^W,
 \qquad\operatorname{sgn}\frac{n_j}{T_j}=\operatorname{sgn}A.
\]
Let $\delta_i=m_i-X_i+(A-n_j)Y_i$. Direct expansion gives
\begin{equation}\label{eq:erroridentity}
 E_j=\sum_{i\in j}\delta_i^2+T_j\sum_{i\in j}Y_i^2.
\end{equation}
Thus \eqref{eq:orientedintegrality} is equivalent to integrality of
$\sum_j n_j\sum_i\delta_i^2/T_j$.

We record why signed indices cause no gap. In the quadratic integer
ring with $s^2=A^2-1$, send $s$ to $n_j-A$ modulo $T_j$. This is a
ring homomorphism because
\[
 (n_j-A)^2-(A^2-1)=n_j^2-2An_j+1=-T_j.
\]
The unit $A+s$ maps to $n_j$, whose inverse is $2A-n_j$, since
$n_j(2A-n_j)=1\pmod{T_j}$. Therefore for every signed integer $k_i$,
\begin{equation}\label{eq:signedcongruence}
 X_i+(n_j-A)Y_i\equiv n_j^{k_i}\pmod{T_j}.
\end{equation}
Here negative powers on the right mean powers of this modular unit.

Put $b=\lfloor\sqrt{W-1}\rfloor$, so $b\ge16$. Let $J$ be the set
of nonzero-base blocks, $s=|J|\le5$, and put
\[
 C=\prod_{j\in J}n_j^b\ne0.
\]
The integer $C$ appears only in the proof. For a request in an active
block define
\[
 d_i=Cm_i-\frac{C}{n_j^b}n_j^{b+k_i}\in\Z.
\]
The exponents $b,b+k_i$ are nonnegative. Equation
\eqref{eq:signedcongruence} implies
$C\delta_i\equiv d_i\pmod{T_j}$, and hence
$C^2\delta_i^2-d_i^2$ is a multiple of $T_j$.
Multiplying the integral sum by $C^2$ now shows that
\[
 \Lambda=\sum_{j\in J}\frac{n_j\sum_{i\in j}d_i^2}{T_j}
\]
is an integer. This step clears each negative-index denominator before
any estimate is used.

Every relevant absolute value is at most $b$, so
\[
 |d_i|\le b^{sb+1}+b^{(s+1)b}\le2b^{(s+1)b}.
\]
Writing $r\le64$ for the total number of requests, we obtain
\begin{align*}
 \sum_{j\in J}|n_j|\sum_{i\in j}d_i^2
 &\le4r\,b^{2(s+1)b+1}\\
 &\le4r\,b^{12b+1}<W^{6b+2}\le W^W.
\end{align*}
The strict inequality follows from $4r\le256<W$ and $b^2<W$.
The last inequality uses $6b+2\le b^2+1\le W$, valid for $b\ge16$.
Since every denominator has absolute value greater than $W^W$,
$|\Lambda|<1$. All its nonzero summands have sign $\operatorname{sgn}A$.
Thus the integer $\Lambda$ vanishes, and every $d_i=0$. Dividing by
$C$ gives $m_i=n_j^{k_i}$ in $\Q$. If $|n_j|\ge2$ and $k_i<0$,
the integer equality $m_i n_j^{-k_i}=1$ is impossible.

Conversely, exact powers imply $T_j\mid\delta_i$ in every active
block. Equation~\eqref{eq:erroridentity} makes $E_j/T_j$ integral
there. Every inactive block contributes zero, proving
\eqref{eq:orientedintegrality}.
\end{proof}

\subsection{Clearing the denominators without adding a primitive}

Introduce one unrestricted integer $z$ and define the ordinary polynomial
\begin{equation}\label{eq:fullpool}
 \mathcal R=
 \sum_{j=0}^{4}n_jE_j\!\prod_{\substack{0\le h\le4\\h\ne j}}T_h
 -z\prod_{h=0}^{4}T_h.
\end{equation}
Since every $T_h\ne0$ under the retained frame, $\mathcal R=0$ is
equivalent to \eqref{eq:orientedintegrality}. The denominator notation
has disappeared from the actual equation. Existing multiplication wires
$Z_i=AY_i$ allow the completely literal definition
\[
 E_j=\sum_{i\in j}
 \left((m_i-D_i)^2-Y_i^2
       -2n_jY_i(m_i-D_i-Z_i)\right).
\]
Thus $E_j$ has degree at most three and $T_j$ degree two. Formula
\eqref{eq:fullpool} has degree at most twelve. Its emitted coefficient
list has a nonzero degree-12 part, so its square has exact degree 24.

For the two existing fixed-base pools, the old private variables
$K_j,L_j$ satisfy
\[
 K_j-2H_j-(2A-n_j)L_j=0,\qquad S_j-n_jK_j+L_j=0.
\]
Adding $n_j$ times the first row to the second gives $E_j-T_jL_j$.
The construction verifies this identity directly against the literal
parent rows; the block membership is not inferred from descriptive
variable names.

\subsection{Activating every block without assuming a positive radix}

The five actual bases are
\[
 t,\quad B,\quad BH,\quad q,\quad q+1,
\]
with respectively $1,1,1,3,10$ requests. All their canonical signed
Pell coordinates and the frame bounds are recovered before any deleted
power row is used. The unchanged source supplies
\begin{gather*}
 t=128+4h^2,\qquad K=t^4,\qquad F\text{ requested as }t^K,\\
 C=2c,\qquad B\mu=C+F+1,\qquad
 H=2+4\sum(\text{actual CRT representatives}),\\
 q\text{ requested as }(BH)^x.
\end{gather*}
The following order is essential.
\begin{enumerate}
\item Since $t\ge128$, its block is active. Theorem~\ref{thm:pool}
gives $F=t^K$. In particular $F$ is even.
\item The literal $C=2c$ is even. Hence $C+F+1$ is odd and nonzero.
The retained equality $B\mu=C+F+1$ forces $B$ odd and nonzero, without
using the sign of $C$.
\item $H\equiv2\pmod4$, so $H\ne0$ and $|H|\ge2$. The active
$BH$ block gives $x\ge0$ and $q=(BH)^x$.
\item If $x=0$, then $q=1$. If $x\ge1$, then $q$ is nonzero and
even. Thus $q\notin\{0,-1\}$, and both remaining bases are nonzero.
\end{enumerate}
All five blocks are now active and all original power assertions hold.
The proof has not changed a CRT representative, parity condition,
radix constructor, or shared residual to obtain this conclusion.

For a new zero, restore each singleton modulus as $T_j$ and its old
quotient as $-\delta_j/T_j$. For a fixed-base block restore
\[
 L_j=E_j/T_j,\qquad K_j=2H_j+(2A-n_j)L_j.
\]
These are integers by the recovered powers. Every retained coordinate
is unchanged. Conversely, at an old zero the old singleton quotient
$Q_j$ gives $\delta_j=-T_jQ_j$, so
$E_j/T_j=T_jQ_j^2+Y_j^2$. The explicit polynomial choice
\[
 z=\sum_{j=0}^{2}n_j(T_jQ_j^2+Y_j^2)
       +qL_q+(q+1)L_{q+1}
\]
satisfies the new certificate. Thus the transformation is pointwise
existential equivalence over the shared source coordinates.

\subsection{A degree-16 companion}

Use oriented pooling only on the first three bases, with quotient $z_0$:
\[
 \mathcal R_0=
 \sum_{j=0}^{2}n_jE_j\!\prod_{\substack{0\le h\le2\\h\ne j}}T_h
 -z_0\prod_{h=0}^{2}T_h.
\]
The preceding parity argument already makes $q,q+1$ nonzero. Two
nonzero consecutive integers have the same sign. Their denominators
therefore have the same sign, and the unweighted two-base certificate
from package 168 applies \cite{source168}:
\begin{equation}\label{eq:splitpool}
 \mathcal R_1=E_3T_4+E_4T_3-z_1T_3T_4.
\end{equation}
For completeness from an old zero take
$z_0=\sum_{j=0}^{2}n_j(T_jQ_j^2+Y_j^2)$ and
$z_1=L_q+L_{q+1}$. For soundness the same bounded-error proof uses
$C=(q(q+1))^b$ and the estimate
\[
 4\cdot13\,b^{6b}<W^{3b+1}\le W^W.
\]
There are no opposite-sign errors to cancel. The first new residual has
exact degree eight, the second exact degree five. Their squares and the
144 retained quadratic squares give exact degree sixteen.

\subsection{Final variable counts and exact projections}

In the parent 221-coordinate chart the removed private coordinates are
\[
 145,146,147,148,149,150,215,216,217,218,
\]
and the removed rows are
\[
 102,103,104,105,106,107,153,154,155,156.
\]
The checker verifies that none of these coordinates occurs in a
retained row. The one-certificate source has $221-10+1=212$ variables
and 151 rows; the split source has 213 variables and 152 rows.

The same five integral projections along directions in the common
kernel of the quadratic Hessians, followed by the same affine
consequence, remove six syntax coordinates. At every step the actual
pivot slope is $1$ or $-1$, the restoration has integer coefficients,
and the new pooling expressions are invariant along the projected
directions. The independent checker expands the shears directly rather
than using the producer's directional-derivative method. Hence the final
counts are 206 and 207.

Both equations share exactly 205 coordinates $x_0,\ldots,x_{204}$,
the same 144 quadratic residuals and the same sixteen request tuples.
For $P_{206}$ the only new coordinate is $z=x_{205}$; for $P_{207}$
they are $z_0=x_{205},z_1=x_{206}$. Appendix B specifies all those
residuals and request tuples literally, so its two short formulas define
both complete integer polynomials without importing any unspecified
universal polynomial.


\section{Logical calibration and formalization boundary}

The two source families directly encode the C222 numerical calculus;
neither construction simulates K215 or instantiates Jones's universal
polynomial. The inherited chain is
\begin{align*}
 \text{integer polynomial zero}
 &\longleftrightarrow\text{finite square-shell C222 refutation}\\
 &\longleftrightarrow\text{original C222 derivation of code 1}\\
 &\longleftrightarrow\text{finite ZF refutation}\\
 &\longleftrightarrow\text{finite ZFC refutation}.
\end{align*}
The square-shell syntax change is an effective recoding, not an assertion
that its numerical constructors literally equal the older Lean ones.
The two logical interfaces are
\begin{center}\small\ttfamily
KEndpoint.c222\_iff\_zf\_refutation\\
ConstructibleChoice.zf\_refutation\_iff\_zfc\_refutation.
\end{center}
Their inspected source files are included \cite{endpoint,choice}.
The numerical-calculus endpoint has no Turing-machine argument, even
though its source file also contains machine-specific theorems.

The full C222/ZF calibration is a substantial inherited theorem. The
lower direction uses a model construction through a tagged domain and
an internally well-founded quotient component; it is not the invalid
assertion that syntactic consistency yields a transitive set model.
The ZF/ZFC step is the constructible-model equiconsistency argument.
The archive identifies these dependencies but does not recompile their
entire transitive Lean import closure. Inspecting two endpoint proof
terms alone would not establish every imported theorem's assumptions.

On the arithmetic side, the inherited nontrivial inputs include signed
Pell classification and canonical-index recovery, the common frame and
orientation bound, generalized binomial CRT coding, the local
proof-calculus compiler, and large-evaluation multiplicity recovery.
The new proofs state their exact interfaces and preserve or strengthen
their hypotheses. The pooled replacement is even pointwise
existentially equivalent over all retained source coordinates.

All new proof operations are finite arithmetic, induction, polynomial
identities, finite groups/rings, elementary lifting, and CRT. They can
be carried out internally in ZFC and do not quantify only over
externally standard witnesses. In a nonstandard model the same internal
lemmas relate its coded polynomial tuples and coded finite derivations.
This is the intended meaning of \eqref{eq:target}.

\section{Reproducibility and evidence}

\subsection{What is completely specified}

The archive includes the LaTeX article and compiled PDF; all four
polynomial residual lists; three full collected polynomial expansions;
the final-coordinate pool blocks and their factored definitions; exact
source changes; integral projection certificates; selected inherited
proofs and endpoint source files; independent checkers; an evaluator;
and checksums. A finite residual list is already a complete polynomial
definition, including for $P_{206}$ whose full squared expansion is not
materialized.

\subsection{Checks and their proper interpretation}

The center endpoint checker imports no producer algebra. It compares
every residual by variable names, checks \eqref{eq:globalidentity}, and
independently collects every coefficient of both quartics. Six deliberate
mutations are rejected. A separate checker authenticates the 21 frame
atoms and the exact parity calculation. Two independently implemented
finite-ring checks cover 59,880 and 2,040 center congruences; the larger
run includes 2,351 nontrivial local order-eight cases. Both mathematical
negative controls are retained.

The pooling audit reconstructs the five blocks from the parent rows,
checks privacy of all ten deleted coordinates, independently replays
the five radical-direction projections and the affine consequence,
and compares every final residual coefficient. Its unbounded proof was
separately scrutinized, including negative indices, negative centers,
unit bases, temporarily zero bases, and the order of the parity
bootstrap. Complete finite fixtures supplement, rather than replace,
those quantified arguments.

The standard-library entry point is
\begin{verbatim}
python3 -B verify.py
python3 -B -O verify.py
\end{verbatim}
It authenticates the release files when a manifest is present, runs the
finite checks, and reconstructs the new equations in temporary
directories. The included receipts record the actual completed runs.
Disabling Python assertions does not remove the substantive checks:
explicit exceptions implement failures.

The evaluator accepts a literal integer tuple of the correct length and
computes the chosen polynomial. It is not a root search or a decision
procedure for consistency. Rebuild commands and relative paths are in
\texttt{README.md}; no network is needed for the delivered checks.

\section{Research implications and next concrete targets}

The first result changes the question asked of a free parameter. The
center need not accept every nonnegative slack at a fixed frame; the
frame itself can be selected from a cofinal family. Removing the odd
part of a finite unit order and using the even frame forces a square
class that one scaled square can realize. This suggests examining other
free Diophantine parameters whose completeness proofs currently use a
universal three- or four-square slack. Any further application must
prove its own residue-class and parity conditions; Theorem~\ref{thm:center}
does not permit indiscriminate replacement of nonnegativity witnesses.

The second result makes reciprocal certificates usable before all
semantic signs are known. The weight $n_j$ aligns denominator signs,
and allowing inactive zero bases separates the recognition lemma from
the later nonvanishing proof. The latter is particularly cheap in this
source because an even power and an even CRT index already make a
retained product odd. A useful next optimization would lower the degree
of the five-denominator certificate while preserving this sign and
activation argument. Eliminating the orienting factors without a new
proof would reintroduce cancellation.

Neither reduction settles the variable/degree/height optimum. The
37-variable large-coefficient quartic and the nine/ten-variable
high-degree branches remain separate competitors. For integration,
the safest comparison is the literal matched-parent table and the
exact source chart used for pooling. An end-to-end formalization of
the arithmetic compiler, including the two new lemmas, is a meaningful
further target; it is not represented here as already complete.

\appendix
\section{Complete balanced quartic \texorpdfstring{$P_{229}$}{P229}}

All $x_0,\ldots,x_{228}$ range independently over $\Z$. The equation is
\[
 P_{229}(x_0,\ldots,x_{228})=\sum_{i=0}^{175}R_i^2=0.
\]
The following is the entire literal residual list. Variable names in
the companion TSV file are explanatory metadata and impose no domain
restrictions.
% Generated from equation.json by code/build_one_square.py
\begingroup\small
\begin{equation*}\begin{split}
R_{0}={}&x_{3} - x_{4} - 2x_{4} x_{227}
\end{split}\end{equation*}
\begin{equation*}\begin{split}
R_{1}={}&x_{0} - x_{0} x_{2} - x_{2} + x_{7}
\end{split}\end{equation*}
\begin{equation*}\begin{split}
R_{2}={}&-1 - x_{2} x_{5} + x_{6} x_{7} - 2x_{227}
\end{split}\end{equation*}
\begin{equation*}\begin{split}
R_{3}={}&-1 + x_{1} - x_{6} x_{7}
\end{split}\end{equation*}
\begin{equation*}\begin{split}
R_{4}={}&1 + x_{4} + x_{16} - x_{18} - 2x_{18} x_{227}
\end{split}\end{equation*}
\begin{equation*}\begin{split}
R_{5}={}&-1 - x_{8}^{2} - 2x_{8} x_{9} - x_{9}^{2} + x_{19}
\end{split}\end{equation*}
\begin{equation*}\begin{split}
R_{6}={}&-10x_{8} + x_{8}^{2} - 6x_{8} x_{9} - 8x_{8} x_{19} - 20x_{9} - 11x_{9}^{2} - 16x_{9} x_{19}\\
&\quad -15x_{19} - 4x_{19}^{2} + x_{20}
\end{split}\end{equation*}
\begin{equation*}\begin{split}
R_{7}={}&-2 + 2x_{8} - x_{8}^{2} - 4x_{8} x_{20} - 4x_{20} - 4x_{20}^{2} + x_{21}
\end{split}\end{equation*}
\begin{equation*}\begin{split}
R_{8}={}&1 + 6x_{8} - 36x_{8}^{2} - 24x_{8} x_{21} + 2x_{21} - 4x_{21}^{2} + x_{22}
\end{split}\end{equation*}
\begin{equation*}\begin{split}
R_{9}={}&-2 - x_{9} - x_{9}^{2} - 2x_{9} x_{10} - x_{10} - x_{10}^{2} + x_{23}
\end{split}\end{equation*}
\begin{equation*}\begin{split}
R_{10}={}&-8 + 7x_{8} - x_{8}^{2} + 4x_{8} x_{10} - 4x_{8} x_{23} - 4x_{10} + 4x_{23}\\
&\quad + x_{24} - x_{46}
\end{split}\end{equation*}
\begin{equation*}\begin{split}
R_{11}={}&-1 + x_{8} - x_{8}^{2} - 2x_{8} x_{10} + x_{10} - x_{10}^{2} + x_{25}
\end{split}\end{equation*}
\begin{equation*}\begin{split}
R_{12}={}&23 + 16x_{8} - 8x_{8} x_{9} - 8x_{8} x_{10} - 16x_{8} x_{19} - 8x_{8} x_{25} + 19x_{9}\\
&\quad + 9x_{9}^{2} - 36x_{9} x_{10} - 12x_{9} x_{25} + 16x_{10} + 8x_{10}^{2} - 24x_{10} x_{19} - 8x_{19} x_{25}\\
&\quad -x_{20} + x_{26} - x_{31}
\end{split}\end{equation*}
\begin{equation*}\begin{split}
R_{13}={}&-4x_{24}^{2} - 8x_{24} x_{26} - 4x_{26}^{2} + x_{27}
\end{split}\end{equation*}
\begin{equation*}\begin{split}
R_{14}={}&-2 - 8x_{8} + 8x_{8} x_{23} + 8x_{8} x_{25} - 8x_{9} + 8x_{9} x_{23} + 8x_{9} x_{25}\\
&\quad -4x_{19} + 8x_{23} - 4x_{23}^{2} - 8x_{23} x_{25} + 8x_{25} - 4x_{25}^{2} + x_{28}
\end{split}\end{equation*}
\begin{equation*}\begin{split}
R_{15}={}&6 + 4x_{8} + 32x_{8} x_{9} + 16x_{8} x_{19} + 16x_{8} x_{25} + 8x_{8} x_{28} + 24x_{9}\\
&\quad -32x_{9} x_{25} - 16x_{9} x_{28} + 12x_{19} - 16x_{19} x_{25} - 8x_{19} x_{28} - x_{20} + 8x_{25}\\
&\quad -16x_{25}^{2} - 16x_{25} x_{28} + 4x_{28} - 4x_{28}^{2} + x_{29}
\end{split}\end{equation*}
\begin{equation*}\begin{split}
R_{16}={}&22 - 5x_{8} - x_{8}^{2} - 8x_{8} x_{9} - 12x_{8} x_{10} - 4x_{8} x_{23} - 8x_{9}\\
&\quad -16x_{9} x_{10} - 16x_{9} x_{23} - 20x_{10} - 20x_{10}^{2} - 24x_{10} x_{23} - 28x_{23} - 4x_{23}^{2}\\
&\quad + x_{30}
\end{split}\end{equation*}
\begin{equation*}\begin{split}
R_{17}={}&19 - 25x_{8} - 7x_{8}^{2} - 8x_{8} x_{9} - 30x_{8} x_{10} - 16x_{8} x_{25} + 3x_{9}\\
&\quad + 5x_{9}^{2} - 4x_{9} x_{25} - 27x_{10} - 21x_{10}^{2} - 24x_{10} x_{25} - 6x_{23} - 17x_{25}\\
&\quad -4x_{25}^{2} + x_{31}
\end{split}\end{equation*}
\begin{equation*}\begin{split}
R_{18}={}&2 - 4x_{30}^{2} - 8x_{30} x_{31} - 4x_{31}^{2} + x_{32}
\end{split}\end{equation*}
\begin{equation*}\begin{split}
R_{19}={}&-19x_{8} + 8x_{8}^{2} + 8x_{8} x_{19} + 8x_{8} x_{23} - 46x_{9} - 24x_{9}^{2} + 24x_{9} x_{19}\\
&\quad + 24x_{9} x_{23} - 2x_{10} + 4x_{19} - 4x_{19}^{2} - 8x_{19} x_{23} + 14x_{23} - 4x_{23}^{2}\\
&\quad -x_{25} + x_{33}
\end{split}\end{equation*}
\begin{equation*}\begin{split}
R_{20}={}&-3 - 6x_{8} - 4x_{10} - 2x_{25} - 4x_{33} - 4x_{33}^{2} + x_{34}
\end{split}\end{equation*}
\begin{equation*}\begin{split}
R_{21}={}&8 + 42x_{8} + 7x_{8}^{2} + 28x_{8} x_{20} + 48x_{9} - 32x_{9}^{2} - 32x_{9} x_{19}\\
&\quad + 64x_{9} x_{20} - 40x_{19}^{2} + 32x_{19} x_{20} + 14x_{20} - 4x_{20}^{2} + x_{35} + 2x_{38}
\end{split}\end{equation*}
\begin{equation*}\begin{split}
R_{22}={}&1 + 6x_{8} - 36x_{8}^{2} - 24x_{8} x_{35} + 2x_{35} - 4x_{35}^{2} + x_{36}
\end{split}\end{equation*}
\begin{equation*}\begin{split}
R_{23}={}&-x_{9}^{2} - 4x_{9} x_{36} - 4x_{36}^{2} + x_{37}
\end{split}\end{equation*}
\begin{equation*}\begin{split}
R_{24}={}&-6 - 2x_{8} - 4x_{9} - 14x_{19} - 16x_{19}^{2} + x_{38}
\end{split}\end{equation*}
\begin{equation*}\begin{split}
R_{25}={}&24 + 21x_{9} - 21x_{9}^{2} - 18x_{9} x_{10} - 24x_{9} x_{23} - 24x_{9} x_{38} + 3x_{10}\\
&\quad -9x_{10}^{2} - 16x_{10} x_{38} + 3x_{23} - 12x_{23}^{2} - 8x_{23} x_{38} + 4x_{38} - 4x_{38}^{2}\\
&\quad + x_{39} + 2x_{46}
\end{split}\end{equation*}
\begin{equation*}\begin{split}
R_{26}={}&16 - 10x_{8} - 20x_{8}^{2} - 16x_{8} x_{10} - 24x_{8} x_{25} - 24x_{8} x_{39} - 12x_{10}\\
&\quad -16x_{10} x_{25} - 16x_{10} x_{39} - 14x_{25} - 4x_{25}^{2} - 8x_{25} x_{39} + 4x_{39} - 4x_{39}^{2}\\
&\quad + x_{40}
\end{split}\end{equation*}
\begin{equation*}\begin{split}
R_{27}={}&-x_{10}^{2} - 2x_{10} x_{11} - x_{11}^{2} + x_{41}
\end{split}\end{equation*}
\begin{equation*}\begin{split}
R_{28}={}&2 + x_{8}^{2} - 4x_{8} x_{41} + 2x_{10} - 2x_{10}^{2} - 8x_{10} x_{41} - 2x_{25}\\
&\quad + 4x_{41} - 4x_{41}^{2} + x_{42}
\end{split}\end{equation*}
\begin{equation*}\begin{split}
R_{29}={}&-31 - 16x_{10} - 16x_{10} x_{23} + 16x_{10} x_{25} - 8x_{23} + 12x_{23}^{2} - 8x_{23} x_{25}\\
&\quad + 24x_{25} - 4x_{25}^{2} + x_{43} - 4x_{46}
\end{split}\end{equation*}
\begin{equation*}\begin{split}
R_{30}={}&-4x_{10} - 16x_{25}^{2} - 32x_{25} x_{41} - 16x_{25} x_{42} - 16x_{25} x_{43} - 4x_{41} - 16x_{41}^{2}\\
&\quad -16x_{41} x_{42} - 16x_{41} x_{43} - 2x_{42} - 4x_{42}^{2} - 8x_{42} x_{43} - 4x_{43}^{2} + x_{44}
\end{split}\end{equation*}
\begin{equation*}\begin{split}
R_{31}={}&-16x_{8}^{2} + 8x_{8} x_{10} + 32x_{8} x_{25} + 16x_{8} x_{28} + x_{10} - x_{10}^{2} - 8x_{10} x_{25}\\
&\quad -4x_{10} x_{28} - 16x_{25}^{2} - 16x_{25} x_{28} - 4x_{28}^{2} + x_{45}
\end{split}\end{equation*}
\begin{equation*}\begin{split}
R_{32}={}&-1 - 8x_{10} - 4x_{10}^{2} + 8x_{10} x_{23} + 8x_{23} - 4x_{23}^{2} + x_{46}
\end{split}\end{equation*}
\begin{equation*}\begin{split}
R_{33}={}&-1 - 9x_{10}^{2} + 24x_{10} x_{23} + 12x_{10} x_{46} - 16x_{23}^{2} - 16x_{23} x_{46} - 4x_{46}^{2}\\
&\quad + x_{47}
\end{split}\end{equation*}
\begin{equation*}\begin{split}
R_{34}={}&1 + x_{9} - x_{9}^{2} - 4x_{9} x_{10} - 4x_{9} x_{47} - 4x_{10}^{2} - 8x_{10} x_{47}\\
&\quad -4x_{47}^{2} + x_{48}
\end{split}\end{equation*}
\begin{equation*}\begin{split}
R_{35}={}&-1 - 4x_{48} - 4x_{48}^{2} + x_{49}
\end{split}\end{equation*}
\begin{equation*}\begin{split}
R_{36}={}&-1 - x_{10}^{2} - 8x_{10} x_{48} - 4x_{10} x_{49} - 16x_{48}^{2} - 16x_{48} x_{49} - 4x_{49}^{2}\\
&\quad + x_{50}
\end{split}\end{equation*}
\begin{equation*}\begin{split}
R_{37}={}&2 + x_{9} - x_{9}^{2} - 4x_{9} x_{10} - 4x_{9} x_{50} - 4x_{10}^{2} - 8x_{10} x_{50}\\
&\quad -4x_{50}^{2} + x_{51}
\end{split}\end{equation*}
\begin{equation*}\begin{split}
R_{38}={}&-4x_{45}^{2} - 8x_{45} x_{51} - 4x_{51}^{2} + x_{52}
\end{split}\end{equation*}
\begin{equation*}\begin{split}
R_{39}={}&-9 + 7x_{8} - x_{8}^{2} - 8x_{8} x_{45} - 4x_{8} x_{52} + 24x_{45} - 16x_{45}^{2}\\
&\quad -16x_{45} x_{52} + 12x_{52} - 4x_{52}^{2} + x_{53}
\end{split}\end{equation*}
\begin{equation*}\begin{split}
R_{40}={}&-1 - 4x_{53} - 4x_{53}^{2} + x_{54}
\end{split}\end{equation*}
\begin{equation*}\begin{split}
R_{41}={}&-x_{10}^{2} - 8x_{10} x_{53} - 4x_{10} x_{54} - 16x_{53}^{2} - 16x_{53} x_{54} - 4x_{54}^{2} + x_{55}
\end{split}\end{equation*}
\begin{equation*}\begin{split}
R_{42}={}&-x_{8} - x_{8}^{2} - 4x_{8} x_{10} - 4x_{8} x_{55} - 4x_{10} - 4x_{10}^{2} - 8x_{10} x_{55}\\
&\quad -4x_{55} - 4x_{55}^{2} + x_{56}
\end{split}\end{equation*}
\begin{equation*}\begin{split}
R_{43}={}&-10 - 12x_{56} - 4x_{56}^{2} + x_{57}
\end{split}\end{equation*}
\begin{equation*}\begin{split}
R_{44}={}&2 + 4x_{10} + 16x_{10} x_{56} + 8x_{10} x_{57} - 4x_{23} - 16x_{23} x_{56} - 8x_{23} x_{57}\\
&\quad -x_{46} + 8x_{56} - 16x_{56}^{2} - 16x_{56} x_{57} + 4x_{57} - 4x_{57}^{2} + x_{58}
\end{split}\end{equation*}
\begin{equation*}\begin{split}
R_{45}={}&x_{9} - x_{9}^{2} + 8x_{9} x_{10} - 8x_{9} x_{23} - 4x_{9} x_{58} - 16x_{10}^{2} + 32x_{10} x_{23}\\
&\quad + 16x_{10} x_{58} - 16x_{23}^{2} - 16x_{23} x_{58} - 4x_{58}^{2} + x_{59}
\end{split}\end{equation*}
\begin{equation*}\begin{split}
R_{46}={}&-10 - 12x_{51} - 4x_{51}^{2} + x_{60}
\end{split}\end{equation*}
\begin{equation*}\begin{split}
R_{47}={}&1 - 16x_{51}^{2} - 16x_{51} x_{59} - 16x_{51} x_{60} - 4x_{59}^{2} - 8x_{59} x_{60} - 4x_{60}^{2}\\
&\quad + x_{61}
\end{split}\end{equation*}
\begin{equation*}\begin{split}
R_{48}={}&x_{10} - x_{10}^{2} - 8x_{10} x_{59} - 4x_{10} x_{61} - 16x_{59}^{2} - 16x_{59} x_{61} - 4x_{61}^{2}\\
&\quad + x_{62}
\end{split}\end{equation*}
\begin{equation*}\begin{split}
R_{49}={}&-6x_{8} - 3x_{8}^{2} - 18x_{8} x_{9} - 16x_{9} - 27x_{9}^{2} + 4x_{9} x_{10} + 3x_{10}\\
&\quad + 3x_{10}^{2} - 4x_{10} x_{19} - 19x_{19} - 8x_{19}^{2} + x_{20} + x_{63} - 4x_{76}
\end{split}\end{equation*}
\begin{equation*}\begin{split}
R_{50}={}&1 - 8x_{8} - 4x_{8} x_{10} + 8x_{8} x_{63} - 16x_{9} - 8x_{9} x_{10} + 16x_{9} x_{63}\\
&\quad + x_{10} + x_{10}^{2} + 4x_{10} x_{19} + 8x_{19} - 8x_{19} x_{63} + 3x_{63} - 4x_{63}^{2}\\
&\quad + x_{64}
\end{split}\end{equation*}
\begin{equation*}\begin{split}
R_{51}={}&-2 - 12x_{8} - 16x_{8} x_{10} + 16x_{8} x_{62} + 16x_{8} x_{64} - 24x_{9} - 32x_{9} x_{10}\\
&\quad + 32x_{9} x_{62} + 32x_{9} x_{64} + 4x_{10} + 4x_{10}^{2} + 16x_{10} x_{19} + 12x_{19} - 16x_{19} x_{62}\\
&\quad -16x_{19} x_{64} - 4x_{62}^{2} - 8x_{62} x_{64} - 4x_{63} - 2x_{64} - 4x_{64}^{2} + x_{65}
\end{split}\end{equation*}
\begin{equation*}\begin{split}
R_{52}={}&-2 - 4x_{65} - 4x_{65}^{2} + x_{66}
\end{split}\end{equation*}
\begin{equation*}\begin{split}
R_{53}={}&-x_{13} - 2x_{13} x_{227} + x_{67}
\end{split}\end{equation*}
\begin{equation*}\begin{split}
R_{54}={}&x_{12} - x_{13} - 2x_{22} x_{67} - x_{68}
\end{split}\end{equation*}
\begin{equation*}\begin{split}
R_{55}={}&x_{12} - x_{13} - 4x_{24} x_{67} - 2x_{27} x_{67} - x_{69}
\end{split}\end{equation*}
\begin{equation*}\begin{split}
R_{56}={}&-2x_{29} x_{67} + 2x_{44} x_{67} - 2x_{67} - x_{70} + x_{75}
\end{split}\end{equation*}
\begin{equation*}\begin{split}
R_{57}={}&x_{12} - x_{13} - 4x_{30} x_{67} - 2x_{32} x_{67} - x_{71}
\end{split}\end{equation*}
\begin{equation*}\begin{split}
R_{58}={}&x_{12} - x_{13} - 2x_{34} x_{67} + 6x_{67} - x_{72}
\end{split}\end{equation*}
\begin{equation*}\begin{split}
R_{59}={}&x_{12} - x_{13} - 4x_{36} x_{67} - 2x_{37} x_{67} - x_{73}
\end{split}\end{equation*}
\begin{equation*}\begin{split}
R_{60}={}&x_{12} - x_{13} - 2x_{40} x_{67} - x_{74}
\end{split}\end{equation*}
\begin{equation*}\begin{split}
R_{61}={}&x_{12} - x_{13} - 2x_{44} x_{67} - 4x_{67} - x_{75}
\end{split}\end{equation*}
\begin{equation*}\begin{split}
R_{62}={}&-x_{8} x_{9} + x_{8} x_{10} + x_{9} + x_{9} x_{10} - x_{10} - x_{10}^{2} + x_{76}
\end{split}\end{equation*}
\begin{equation*}\begin{split}
R_{63}={}&1 + 2x_{11} + x_{11}^{2} - x_{76}^{2} + x_{77}
\end{split}\end{equation*}
\begin{equation*}\begin{split}
R_{64}={}&-x_{77}^{2} + x_{78}
\end{split}\end{equation*}
\begin{equation*}\begin{split}
R_{65}={}&-4x_{65} x_{67} - 2x_{66} x_{67} + 2x_{67} + x_{79}
\end{split}\end{equation*}
\begin{equation*}\begin{split}
R_{66}={}&-x_{12}^{2} + 2x_{12} x_{13} + 2x_{12} x_{79} - x_{13}^{2} - 2x_{13} x_{79} - x_{79}^{2} + x_{80}
\end{split}\end{equation*}
\begin{equation*}\begin{split}
R_{67}={}&-1 - x_{9} - 2x_{9} x_{227} + x_{81}
\end{split}\end{equation*}
\begin{equation*}\begin{split}
R_{68}={}&-x_{13} x_{81} + x_{83} x_{84} + x_{85}
\end{split}\end{equation*}
\begin{equation*}\begin{split}
R_{69}={}&-x_{8} - 2x_{8} x_{227} + x_{82} + 2x_{227}
\end{split}\end{equation*}
\begin{equation*}\begin{split}
R_{70}={}&1 + 2x_{8} + 4x_{8} x_{227} + 4x_{9} + 8x_{9} x_{227} - 2x_{19} - 4x_{19} x_{227}\\
&\quad + x_{83} + 4x_{227}
\end{split}\end{equation*}
\begin{equation*}\begin{split}
R_{71}={}&-x_{12} x_{82} + x_{84}
\end{split}\end{equation*}
\begin{equation*}\begin{split}
R_{72}={}&x_{12} x_{81} + 6x_{13} - 8x_{13} x_{81} - 4x_{13} x_{82} - x_{13} x_{83} - 12x_{13} x_{227} + x_{86}
\end{split}\end{equation*}
\begin{equation*}\begin{split}
R_{73}={}&-x_{12} x_{68} + x_{13} x_{68} + x_{87}
\end{split}\end{equation*}
\begin{equation*}\begin{split}
R_{74}={}&-x_{69} x_{87} + x_{88}
\end{split}\end{equation*}
\begin{equation*}\begin{split}
R_{75}={}&-x_{70} x_{88} + x_{89}
\end{split}\end{equation*}
\begin{equation*}\begin{split}
R_{76}={}&-x_{71} x_{89} + x_{90}
\end{split}\end{equation*}
\begin{equation*}\begin{split}
R_{77}={}&-x_{72} x_{90} + x_{91}
\end{split}\end{equation*}
\begin{equation*}\begin{split}
R_{78}={}&-x_{73} x_{91} + x_{92}
\end{split}\end{equation*}
\begin{equation*}\begin{split}
R_{79}={}&-x_{74} x_{92} + x_{93}
\end{split}\end{equation*}
\begin{equation*}\begin{split}
R_{80}={}&-x_{75} x_{93} + x_{94}
\end{split}\end{equation*}
\begin{equation*}\begin{split}
R_{81}={}&-x_{85} x_{94} + x_{95}
\end{split}\end{equation*}
\begin{equation*}\begin{split}
R_{82}={}&-x_{86} x_{95} + x_{96}
\end{split}\end{equation*}
\begin{equation*}\begin{split}
R_{83}={}&-x_{96}^{2} + x_{97}
\end{split}\end{equation*}
\begin{equation*}\begin{split}
R_{84}={}&-x_{2} - x_{2} x_{16} + x_{16} + x_{98}
\end{split}\end{equation*}
\begin{equation*}\begin{split}
R_{85}={}&-x_{14} x_{98} + x_{99}
\end{split}\end{equation*}
\begin{equation*}\begin{split}
R_{86}={}&-2x_{1}^{2} + 2x_{1} x_{12} + 2x_{1} x_{13} + 2x_{1} x_{99} + 2x_{1} x_{100} - x_{12}^{2} - 2x_{12} x_{99}\\
&\quad -x_{13}^{2} - 2x_{13} x_{100} + x_{17} x_{102} - x_{78} x_{97} - x_{80} x_{97} - x_{99}^{2} - x_{100}^{2}
\end{split}\end{equation*}
\begin{equation*}\begin{split}
R_{87}={}&-x_{2} x_{15} + x_{15} - x_{15} x_{98} + x_{100}
\end{split}\end{equation*}
\begin{equation*}\begin{split}
R_{88}={}&-x_{8} - x_{9} - 4x_{19} - 2x_{22} - 4x_{24} - 2x_{27} - 2x_{29}\\
&\quad -4x_{30} - 2x_{32} - 2x_{34} - 4x_{36} - 2x_{37} - 2x_{40} - 2x_{44}\\
&\quad -4x_{65} - 2x_{66} + x_{101}
\end{split}\end{equation*}
\begin{equation*}\begin{split}
R_{89}={}&-3x_{105} - 6x_{105} x_{221} + x_{107}
\end{split}\end{equation*}
\begin{equation*}\begin{split}
R_{90}={}&x_{103} x_{212} - x_{106} - x_{107} x_{128} - x_{128} - 2x_{128} x_{227}
\end{split}\end{equation*}
\begin{equation*}\begin{split}
R_{91}={}&x_{108} x_{213} - x_{110} x_{131} - x_{111}
\end{split}\end{equation*}
\begin{equation*}\begin{split}
R_{92}={}&x_{112} x_{213} - x_{113} x_{131} - x_{114}
\end{split}\end{equation*}
\begin{equation*}\begin{split}
R_{93}={}&x_{115} x_{213} - x_{116} x_{131} - x_{117}
\end{split}\end{equation*}
\begin{equation*}\begin{split}
R_{94}={}&x_{118} x_{213} - x_{119} x_{131} - x_{120}
\end{split}\end{equation*}
\begin{equation*}\begin{split}
R_{95}={}&x_{121} x_{214} - x_{123} x_{131} - x_{124}
\end{split}\end{equation*}
\begin{equation*}\begin{split}
R_{96}={}&x_{125} x_{214} - x_{126} x_{131} - x_{127}
\end{split}\end{equation*}
\begin{equation*}\begin{split}
R_{97}={}&-x_{17} x_{128} + x_{131}
\end{split}\end{equation*}
\begin{equation*}\begin{split}
R_{98}={}&-x_{131} x_{135} + x_{134} x_{212} - x_{136}
\end{split}\end{equation*}
\begin{equation*}\begin{split}
R_{99}={}&-x_{131} x_{138} + x_{137} x_{213} - x_{139}
\end{split}\end{equation*}
\begin{equation*}\begin{split}
R_{100}={}&x_{140} - 2x_{140} x_{142} + x_{143}
\end{split}\end{equation*}
\begin{equation*}\begin{split}
R_{101}={}&-x_{140} x_{143} - x_{143} + x_{144}
\end{split}\end{equation*}
\begin{equation*}\begin{split}
R_{102}={}&-1 + x_{143}^{2} - x_{144}^{2} + x_{224}^{2}
\end{split}\end{equation*}
\begin{equation*}\begin{split}
R_{103}={}&-x_{141} x_{146} - 7x_{146} + x_{147}
\end{split}\end{equation*}
\begin{equation*}\begin{split}
R_{104}={}&2x_{145} + x_{145}^{2} + x_{146}^{2} - x_{147}^{2}
\end{split}\end{equation*}
\begin{equation*}\begin{split}
R_{105}={}&1 - x_{4} - 32x_{141} x_{148} + x_{145} + 16x_{146} - x_{147} + 33x_{148}
\end{split}\end{equation*}
\begin{equation*}\begin{split}
R_{106}={}&-12 - 2x_{141} - 4x_{141} x_{227} + x_{149} - 24x_{227} + 4x_{227}^{2}
\end{split}\end{equation*}
\begin{equation*}\begin{split}
R_{107}={}&x_{17} + x_{149} x_{150} - x_{183} - x_{184} - 2x_{184} x_{227} + x_{185}
\end{split}\end{equation*}
\begin{equation*}\begin{split}
R_{108}={}&-1 + 2x_{141} x_{196} - x_{151} + 14x_{196} - x_{196}^{2}
\end{split}\end{equation*}
\begin{equation*}\begin{split}
R_{109}={}&3 + x_{151} x_{152} - x_{183} - x_{184} x_{196} + x_{185} + 6x_{221}
\end{split}\end{equation*}
\begin{equation*}\begin{split}
R_{110}={}&-x_{141} x_{154} - 7x_{154} + x_{155}
\end{split}\end{equation*}
\begin{equation*}\begin{split}
R_{111}={}&2x_{153} + x_{153}^{2} + x_{154}^{2} - x_{155}^{2}
\end{split}\end{equation*}
\begin{equation*}\begin{split}
R_{112}={}&-39 + x_{131} x_{228} - 8x_{141} - 12x_{141} x_{221} - 36x_{221} + 36x_{221}^{2}
\end{split}\end{equation*}
\begin{equation*}\begin{split}
R_{113}={}&1 - x_{103} - x_{141} x_{154} + x_{153} - 3x_{154} + 6x_{154} x_{221}
\end{split}\end{equation*}
\begin{equation*}\begin{split}
R_{114}={}&7 - x_{131} x_{228} + 2x_{141} + x_{156} - 12x_{221}
\end{split}\end{equation*}
\begin{equation*}\begin{split}
R_{115}={}&-1 + x_{104} - x_{145} - 3x_{146} - 6x_{146} x_{221} + x_{147} + x_{156} x_{157}
\end{split}\end{equation*}
\begin{equation*}\begin{split}
R_{116}={}&7 + x_{141} - x_{141} x_{159} - 7x_{159} + x_{160}
\end{split}\end{equation*}
\begin{equation*}\begin{split}
R_{117}={}&1 + 2x_{158} + x_{158}^{2} - 2x_{159} + x_{159}^{2} - x_{160}^{2}
\end{split}\end{equation*}
\begin{equation*}\begin{split}
R_{118}={}&4 - x_{108} + x_{141} - x_{141} x_{159} + x_{158} - 3x_{159} + 6x_{159} x_{221}\\
&\quad -6x_{221}
\end{split}\end{equation*}
\begin{equation*}\begin{split}
R_{119}={}&-3x_{104} - 6x_{104} x_{221} + x_{109}
\end{split}\end{equation*}
\begin{equation*}\begin{split}
R_{120}={}&-x_{141} x_{162} - 7x_{162} + x_{163}
\end{split}\end{equation*}
\begin{equation*}\begin{split}
R_{121}={}&2x_{161} + x_{161}^{2} + x_{162}^{2} - x_{163}^{2}
\end{split}\end{equation*}
\begin{equation*}\begin{split}
R_{122}={}&1 - x_{112} - x_{141} x_{162} + x_{161} - 3x_{162} + 6x_{162} x_{221}
\end{split}\end{equation*}
\begin{equation*}\begin{split}
R_{123}={}&-x_{141} x_{165} - 7x_{165} + x_{166}
\end{split}\end{equation*}
\begin{equation*}\begin{split}
R_{124}={}&2x_{164} + x_{164}^{2} + x_{165}^{2} - x_{166}^{2}
\end{split}\end{equation*}
\begin{equation*}\begin{split}
R_{125}={}&1 - x_{115} - x_{141} x_{165} + x_{164} - 3x_{165} + 6x_{165} x_{221}
\end{split}\end{equation*}
\begin{equation*}\begin{split}
R_{126}={}&-x_{141} x_{168} - 7x_{168} + x_{169}
\end{split}\end{equation*}
\begin{equation*}\begin{split}
R_{127}={}&2x_{167} + x_{167}^{2} + x_{168}^{2} - x_{169}^{2}
\end{split}\end{equation*}
\begin{equation*}\begin{split}
R_{128}={}&1 - x_{118} - x_{141} x_{168} + x_{167} - 3x_{168} + 6x_{168} x_{221}
\end{split}\end{equation*}
\begin{equation*}\begin{split}
R_{129}={}&-x_{141} x_{171} - 7x_{171} + x_{172}
\end{split}\end{equation*}
\begin{equation*}\begin{split}
R_{130}={}&2x_{170} + x_{170}^{2} + x_{171}^{2} - x_{172}^{2}
\end{split}\end{equation*}
\begin{equation*}\begin{split}
R_{131}={}&1 - x_{121} - x_{141} x_{171} + x_{170} - 3x_{171} + 6x_{171} x_{221}
\end{split}\end{equation*}
\begin{equation*}\begin{split}
R_{132}={}&-x_{141} x_{174} - 7x_{174} + x_{175}
\end{split}\end{equation*}
\begin{equation*}\begin{split}
R_{133}={}&2x_{173} + x_{173}^{2} + x_{174}^{2} - x_{175}^{2}
\end{split}\end{equation*}
\begin{equation*}\begin{split}
R_{134}={}&-1 + x_{122} + x_{156} x_{176} - x_{173} - 3x_{174} - 6x_{174} x_{221} + x_{175}
\end{split}\end{equation*}
\begin{equation*}\begin{split}
R_{135}={}&-x_{141} x_{178} - 7x_{178} + x_{179}
\end{split}\end{equation*}
\begin{equation*}\begin{split}
R_{136}={}&2x_{177} + x_{177}^{2} + x_{178}^{2} - x_{179}^{2}
\end{split}\end{equation*}
\begin{equation*}\begin{split}
R_{137}={}&1 - x_{125} - x_{141} x_{178} + x_{177} - 3x_{178} + 6x_{178} x_{221}
\end{split}\end{equation*}
\begin{equation*}\begin{split}
R_{138}={}&-x_{141} x_{181} - 7x_{181} + x_{182}
\end{split}\end{equation*}
\begin{equation*}\begin{split}
R_{139}={}&2x_{180} + x_{180}^{2} + x_{181}^{2} - x_{182}^{2}
\end{split}\end{equation*}
\begin{equation*}\begin{split}
R_{140}={}&1 - x_{129} x_{131} - x_{130} - x_{141} x_{181} + x_{180} - 3x_{181} + 6x_{181} x_{221}
\end{split}\end{equation*}
\begin{equation*}\begin{split}
R_{141}={}&1 - x_{141} x_{184} - 7x_{184} + x_{185}
\end{split}\end{equation*}
\begin{equation*}\begin{split}
R_{142}={}&-1 + 2x_{183} + x_{183}^{2} + x_{184}^{2} - 2x_{185} - x_{185}^{2}
\end{split}\end{equation*}
\begin{equation*}\begin{split}
R_{143}={}&x_{128} + x_{156} x_{186} - x_{183} - 3x_{184} - 6x_{184} x_{221} + x_{185}
\end{split}\end{equation*}
\begin{equation*}\begin{split}
R_{144}={}&-x_{141} x_{188} - 7x_{188} + x_{189}
\end{split}\end{equation*}
\begin{equation*}\begin{split}
R_{145}={}&2x_{187} + x_{187}^{2} + x_{188}^{2} - x_{189}^{2}
\end{split}\end{equation*}
\begin{equation*}\begin{split}
R_{146}={}&1 - x_{131} x_{132} - x_{133} - x_{141} x_{188} + x_{187} - 3x_{188} + 6x_{188} x_{221}
\end{split}\end{equation*}
\begin{equation*}\begin{split}
R_{147}={}&-x_{141} x_{191} - 7x_{191} + x_{192}
\end{split}\end{equation*}
\begin{equation*}\begin{split}
R_{148}={}&2x_{190} + x_{190}^{2} + x_{191}^{2} - x_{192}^{2}
\end{split}\end{equation*}
\begin{equation*}\begin{split}
R_{149}={}&1 - x_{134} - x_{141} x_{191} + x_{190} - 3x_{191} + 6x_{191} x_{221}
\end{split}\end{equation*}
\begin{equation*}\begin{split}
R_{150}={}&-x_{141} x_{194} - 7x_{194} + x_{195}
\end{split}\end{equation*}
\begin{equation*}\begin{split}
R_{151}={}&2x_{193} + x_{193}^{2} + x_{194}^{2} - x_{195}^{2}
\end{split}\end{equation*}
\begin{equation*}\begin{split}
R_{152}={}&1 - x_{137} - x_{141} x_{194} + x_{193} - 3x_{194} + 6x_{194} x_{221}
\end{split}\end{equation*}
\begin{equation*}\begin{split}
R_{153}={}&-13 - 2x_{0}^{2} - 2x_{0} x_{4} - x_{1}^{2} - x_{2}^{2} - x_{4}^{2} + 2x_{8}\\
&\quad + 2x_{8} x_{9} - x_{10}^{2} - x_{11}^{2} - x_{12}^{2} - x_{13}^{2} - x_{14}^{2} - x_{15}^{2}\\
&\quad -x_{16}^{2} - x_{17}^{2} - x_{19} - x_{101}^{2} - x_{104}^{2} - x_{122}^{2} - x_{128}^{2}\\
&\quad + x_{140} - x_{196}^{2} - 36x_{221} - 36x_{221}^{2} - 4x_{227} - 4x_{227}^{2}
\end{split}\end{equation*}
\begin{equation*}\begin{split}
R_{154}={}&x_{196} - x_{226} - 2x_{226} x_{227}
\end{split}\end{equation*}
\begin{equation*}\begin{split}
R_{155}={}&-16 + x_{0} - x_{197}^{2}
\end{split}\end{equation*}
\begin{equation*}\begin{split}
R_{156}={}&-x_{141} x_{199} - 7x_{199} + x_{208}
\end{split}\end{equation*}
\begin{equation*}\begin{split}
R_{157}={}&-1 + x_{198}^{2} + x_{199}^{2} - x_{208}^{2}
\end{split}\end{equation*}
\begin{equation*}\begin{split}
R_{158}={}&-x_{201} - x_{201} x_{202} + x_{209}
\end{split}\end{equation*}
\begin{equation*}\begin{split}
R_{159}={}&-1 + x_{200}^{2} + x_{201}^{2} - x_{209}^{2}
\end{split}\end{equation*}
\begin{equation*}\begin{split}
R_{160}={}&-x_{194}^{2} + x_{210}
\end{split}\end{equation*}
\begin{equation*}\begin{split}
R_{161}={}&x_{199} - 2x_{204} x_{210} - x_{210}
\end{split}\end{equation*}
\begin{equation*}\begin{split}
R_{162}={}&-4x_{199} - 8x_{199} x_{203} + x_{202}
\end{split}\end{equation*}
\begin{equation*}\begin{split}
R_{163}={}&-6 - x_{141} - x_{198} x_{205} + x_{202}
\end{split}\end{equation*}
\begin{equation*}\begin{split}
R_{164}={}&-1 - x_{193} - x_{198} x_{206} + x_{200}
\end{split}\end{equation*}
\begin{equation*}\begin{split}
R_{165}={}&-x_{101} - 4x_{194} x_{207} + x_{201}
\end{split}\end{equation*}
\begin{equation*}\begin{split}
R_{166}={}&-2 + 2x_{0}^{2} + 2x_{0} x_{4} - 2x_{0} x_{146} - 2x_{0} x_{154} + x_{1}^{2} - 2x_{1} x_{184}\\
&\quad + x_{2}^{2} - 2x_{2} x_{174} + x_{4}^{2} - 2x_{4} x_{154} + x_{8}^{2} - 2x_{8} x_{159} + x_{9}^{2}\\
&\quad -2x_{9} x_{162} + x_{10}^{2} - 2x_{10} x_{165} + x_{11}^{2} - 2x_{11} x_{168} + x_{12}^{2} - 2x_{12} x_{171}\\
&\quad + x_{13}^{2} - 2x_{13} x_{178} + x_{14}^{2} - 2x_{14} x_{181} + x_{15}^{2} - 2x_{15} x_{188} + x_{16}^{2}\\
&\quad -2x_{16} x_{191} - 2x_{101} x_{194} - x_{141} x_{211} + x_{145}^{2} + x_{153}^{2} + x_{158}^{2} + 2x_{159}\\
&\quad + x_{161}^{2} + x_{164}^{2} + x_{167}^{2} + x_{170}^{2} + x_{173}^{2} + x_{177}^{2} + x_{180}^{2}\\
&\quad + x_{183}^{2} + x_{187}^{2} + x_{190}^{2} + x_{193}^{2} + 2x_{194}^{2} - 6x_{211} - x_{215}^{2}\\
&\quad -x_{216}^{2} - x_{217} - x_{217}^{2}
\end{split}\end{equation*}
\begin{equation*}\begin{split}
R_{167}={}&-x_{104} x_{212} + x_{128}
\end{split}\end{equation*}
\begin{equation*}\begin{split}
R_{168}={}&-x_{109} x_{213} + x_{128}
\end{split}\end{equation*}
\begin{equation*}\begin{split}
R_{169}={}&-x_{122} x_{214} + x_{128}
\end{split}\end{equation*}
\begin{equation*}\begin{split}
R_{170}={}&-1 + x_{2} - x_{3} x_{222}
\end{split}\end{equation*}
\begin{equation*}\begin{split}
R_{171}={}&-2 - x_{101}^{2} - x_{146}^{2} - x_{154}^{2} + 2x_{159} - x_{159}^{2} - x_{162}^{2}\\
&\quad -x_{165}^{2} - x_{168}^{2} - x_{171}^{2} - x_{174}^{2} - x_{178}^{2} - x_{181}^{2} - x_{184}^{2}\\
&\quad -x_{188}^{2} - x_{191}^{2} + x_{194}^{2} - x_{215}^{2} - x_{216}^{2} - x_{217} - x_{217}^{2}
\end{split}\end{equation*}
\begin{equation*}\begin{split}
R_{172}={}&-1 - 4x_{106}^{2} - 4x_{111}^{2} - 4x_{114}^{2} - 4x_{117}^{2} - 4x_{120}^{2} - 4x_{124}^{2}\\
&\quad -4x_{127}^{2} + x_{128}^{2} - 4x_{130}^{2} - 4x_{133}^{2} - 4x_{136}^{2} - 4x_{139}^{2} - x_{218}^{2}\\
&\quad -x_{219}^{2} - x_{220} - x_{220}^{2}
\end{split}\end{equation*}
\begin{equation*}\begin{split}
R_{173}={}&-1 + x_{222} - 4x_{223} - 4x_{223}^{2}
\end{split}\end{equation*}
\begin{equation*}\begin{split}
R_{174}={}&7 + x_{141} - x_{224}^{2} - 2x_{225}^{2}
\end{split}\end{equation*}
\begin{equation*}\begin{split}
R_{175}={}&-3 - 2x_{0} - 2x_{4} + 2x_{8} - x_{8}^{2} - x_{9}^{2} - x_{10}^{2}\\
&\quad -x_{11}^{2} - x_{12}^{2} - x_{13}^{2} - x_{14}^{2} - x_{15}^{2} - x_{16}^{2} + x_{226}
\end{split}\end{equation*}
\endgroup


\section{Complete definitions of the pooled equations}
All coordinates range independently over the integers. These two
polynomials share the following 205 coordinates and 144 quadratic residuals.
Set
\[
A=x_{132}+1,\qquad
(n_0,n_1,n_2,n_3,n_4)=(x_0,x_3,x_{178},x_{95}-1,x_{95}).
\]
For each row of the following table, use the displayed affine polynomials
as $(m,D,Y,Z)$ in its block. The table is part of the polynomial definition;
Pell notation is unnecessary to evaluate it.
\begin{center}\small
\begin{tabular}{@{}crrrrr@{}}\toprule
Block $j$ & Request & $m$ & $D$ & $Y$ & $Z$\\\midrule
0 & 0 & $x_{2}$ & $x_{175}$ & $x_{176}$ & $x_{177}$\\
\addlinespace[3pt]
1 & 0 & $x_{17}$ & $x_{166}$ & $x_{167}$ & $1 + x_{168}$\\
\addlinespace[3pt]
2 & 0 & $-1 + x_{95}$ & $x_{166}$ & $x_{167}$ & $1 + x_{168}$\\
\addlinespace[3pt]
3 & 0 & $x_{97}$ & $x_{136}$ & $x_{137}$ & $x_{138}$\\
3 & 1 & $x_{114}$ & $x_{157}$ & $x_{158}$ & $x_{159}$\\
3 & 2 & $x_{121}$ & $x_{166}$ & $x_{167}$ & $1 + x_{168}$\\
\addlinespace[3pt]
4 & 0 & $x_{96}$ & $x_{139}$ & $x_{140}$ & $x_{141}$\\
4 & 1 & $x_{101}$ & $x_{142}$ & $-1 + x_{143}$ & $x_{144}$\\
4 & 2 & $x_{104}$ & $x_{145}$ & $x_{146}$ & $x_{147}$\\
4 & 3 & $x_{107}$ & $x_{148}$ & $x_{149}$ & $x_{150}$\\
4 & 4 & $x_{110}$ & $x_{151}$ & $x_{152}$ & $x_{153}$\\
4 & 5 & $x_{113}$ & $x_{154}$ & $x_{155}$ & $x_{156}$\\
4 & 6 & $x_{117}$ & $x_{160}$ & $x_{161}$ & $x_{162}$\\
4 & 7 & $x_{120}$ & $x_{163}$ & $x_{164}$ & $x_{165}$\\
4 & 8 & $x_{125}$ & $x_{169}$ & $x_{170}$ & $x_{171}$\\
4 & 9 & $x_{128}$ & $x_{172}$ & $x_{173}$ & $x_{174}$\\
\bottomrule\end{tabular}\end{center}
Define the following fixed polynomials by finite sums over the table:
\[
T_j=2An_j-n_j^2-1,\qquad
E_j=\sum_{i\in j}\left((m_i-D_i)^2-Y_i^2
                  -2n_jY_i(m_i-D_i-Z_i)\right).
\]
No table entry is an extra unknown. With $Q_0,\ldots,Q_{143}$ given in
full below, the 206-variable equation is
\begin{align*}
P_{206}(x_0,\ldots,x_{205})
&=\sum_{i=0}^{143}Q_i^2+\mathcal R^2=0,\\
\mathcal R
&=\sum_{j=0}^{4}n_jE_j\prod_{\substack{0\le h\le4\\h\ne j}}T_h
  -x_{205}\prod_{h=0}^{4}T_h.
\end{align*}
The 207-variable equation is
\begin{align*}
P_{207}(x_0,\ldots,x_{206})
&=\sum_{i=0}^{143}Q_i^2+\mathcal R_0^2+\mathcal R_1^2=0,\\
\mathcal R_0
&=\sum_{j=0}^{2}n_jE_j\prod_{\substack{0\le h\le2\\h\ne j}}T_h
  -x_{205}T_0T_1T_2,\\
\mathcal R_1&=E_3T_4+E_4T_3-x_{206}T_3T_4.
\end{align*}
These finite formulas, the complete table, and the following list are a
standalone specification of both ordinary integer polynomials. In
particular the symbolic products have the fixed lengths three or five;
there is no variable-length product in a final equation.

% Generated from equation.json by code/build_one_square.py
\begingroup\small
\begin{equation*}\begin{split}
Q_{0}={}&x_{0} - x_{0} x_{4} - x_{4} + x_{7}
\end{split}\end{equation*}
\begin{equation*}\begin{split}
Q_{1}={}&-1 + x_{1} - x_{2} - x_{4} x_{5}
\end{split}\end{equation*}
\begin{equation*}\begin{split}
Q_{2}={}&-1 + x_{1} - x_{6} x_{7}
\end{split}\end{equation*}
\begin{equation*}\begin{split}
Q_{3}={}&1 + x_{2} - x_{3} x_{18} + 2x_{16}
\end{split}\end{equation*}
\begin{equation*}\begin{split}
Q_{4}={}&-1 + 2x_{8} - x_{8}^{2} - 2x_{8} x_{9} + 2x_{9} - x_{9}^{2} + x_{19}
\end{split}\end{equation*}
\begin{equation*}\begin{split}
Q_{5}={}&-2 + 2x_{8} - x_{8}^{2} - 4x_{8} x_{20} + 4x_{20} - 4x_{20}^{2} + x_{21}
\end{split}\end{equation*}
\begin{equation*}\begin{split}
Q_{6}={}&1 + 6x_{8} - 36x_{8}^{2} - 192x_{8} x_{20} - 24x_{8} x_{21} + 16x_{20} - 256x_{20}^{2}\\
&\quad -64x_{20} x_{21} + 2x_{21} - 4x_{21}^{2} + x_{22}
\end{split}\end{equation*}
\begin{equation*}\begin{split}
Q_{7}={}&-2 - x_{9}^{2} - 2x_{9} x_{10} - x_{10}^{2} + x_{23}
\end{split}\end{equation*}
\begin{equation*}\begin{split}
Q_{8}={}&-1 + 2x_{8} - x_{8}^{2} - 2x_{8} x_{10} + 2x_{10} - x_{10}^{2} + x_{24}
\end{split}\end{equation*}
\begin{equation*}\begin{split}
Q_{9}={}&-9 + x_{8} - 145x_{8}^{2} - 196x_{8} x_{9} - 192x_{8} x_{10} - 48x_{8} x_{19} - 4x_{8} x_{23}\\
&\quad -48x_{8} x_{24} + 4x_{9} - 68x_{9}^{2} - 128x_{9} x_{10} - 32x_{9} x_{19} - 8x_{9} x_{23} - 32x_{9} x_{24}\\
&\quad -64x_{10}^{2} - 32x_{10} x_{19} - 32x_{10} x_{24} - 2x_{19} - 4x_{19}^{2} - 8x_{19} x_{24} + 12x_{23}\\
&\quad -4x_{23}^{2} - 4x_{24}^{2} + x_{25}
\end{split}\end{equation*}
\begin{equation*}\begin{split}
Q_{10}={}&-18 + 14x_{8} - 2x_{8}^{2} - 8x_{8} x_{9} - 8x_{8} x_{23} + 24x_{9} - 8x_{9}^{2}\\
&\quad -16x_{9} x_{23} + 24x_{23} - 8x_{23}^{2} - 4x_{25}^{2} + x_{26}
\end{split}\end{equation*}
\begin{equation*}\begin{split}
Q_{11}={}&-6 + 16x_{10} - 16x_{10}^{2} - 16x_{10} x_{23} - 16x_{10} x_{24} + 8x_{23} - 4x_{23}^{2}\\
&\quad -8x_{23} x_{24} + 8x_{24} - 4x_{24}^{2} + x_{27}
\end{split}\end{equation*}
\begin{equation*}\begin{split}
Q_{12}={}&1 + 6x_{8} - 36x_{8}^{2} - 96x_{8} x_{9} - 48x_{8} x_{10} - 24x_{8} x_{19} - 48x_{8} x_{24}\\
&\quad -24x_{8} x_{27} + 8x_{9} - 64x_{9}^{2} - 64x_{9} x_{10} - 32x_{9} x_{19} - 64x_{9} x_{24} - 32x_{9} x_{27}\\
&\quad + 8x_{10} - 16x_{10}^{2} - 16x_{10} x_{19} - 32x_{10} x_{24} - 16x_{10} x_{27} + 2x_{19} - 4x_{19}^{2}\\
&\quad -16x_{19} x_{24} - 8x_{19} x_{27} + 8x_{24} - 16x_{24}^{2} - 16x_{24} x_{27} + 4x_{27} - 4x_{27}^{2}\\
&\quad + x_{28}
\end{split}\end{equation*}
\begin{equation*}\begin{split}
Q_{13}={}&-12 - 29x_{8} - 37x_{8}^{2} - 24x_{8} x_{9} - 112x_{8} x_{10} - 4x_{8} x_{23} - 24x_{8} x_{24}\\
&\quad -39x_{9} - 37x_{9}^{2} - 112x_{9} x_{10} - 24x_{9} x_{23} - 4x_{9} x_{24} - 80x_{10} - 128x_{10}^{2}\\
&\quad -32x_{10} x_{23} - 32x_{10} x_{24} - 12x_{23} - 4x_{23}^{2} - 8x_{24} - 4x_{24}^{2} + x_{29}
\end{split}\end{equation*}
\begin{equation*}\begin{split}
Q_{14}={}&-18 - 10x_{8} - 2x_{8}^{2} - 24x_{8} x_{9} - 32x_{8} x_{10} - 8x_{8} x_{23} - 72x_{9}\\
&\quad -72x_{9}^{2} - 192x_{9} x_{10} - 48x_{9} x_{23} - 96x_{10} - 128x_{10}^{2} - 64x_{10} x_{23} - 24x_{23}\\
&\quad -8x_{23}^{2} - 4x_{29}^{2} + x_{30}
\end{split}\end{equation*}
\begin{equation*}\begin{split}
Q_{15}={}&-12 + 16x_{8} - 4x_{8}^{2} - 8x_{8} x_{10} - 8x_{8} x_{19} - 8x_{8} x_{23} + 16x_{10}\\
&\quad -4x_{10}^{2} - 8x_{10} x_{19} - 8x_{10} x_{23} + 16x_{19} - 4x_{19}^{2} - 8x_{19} x_{23} + 16x_{23}\\
&\quad -4x_{23}^{2} + x_{31}
\end{split}\end{equation*}
\begin{equation*}\begin{split}
Q_{16}={}&-3 - 24x_{8} - 64x_{8}^{2} - 160x_{8} x_{10} - 64x_{8} x_{23} - 32x_{8} x_{24} - 32x_{8} x_{31}\\
&\quad -26x_{10} - 100x_{10}^{2} - 80x_{10} x_{23} - 40x_{10} x_{24} - 40x_{10} x_{31} - 8x_{23} - 16x_{23}^{2}\\
&\quad -16x_{23} x_{24} - 16x_{23} x_{31} - 6x_{24} - 4x_{24}^{2} - 8x_{24} x_{31} - 4x_{31} - 4x_{31}^{2}\\
&\quad + x_{32}
\end{split}\end{equation*}
\begin{equation*}\begin{split}
Q_{17}={}&-5 - 30x_{8} - 36x_{8}^{2} - 96x_{8} x_{9} - 24x_{8} x_{19} - 40x_{9} - 64x_{9}^{2}\\
&\quad -32x_{9} x_{19} - 10x_{19} - 4x_{19}^{2} + x_{20}
\end{split}\end{equation*}
\begin{equation*}\begin{split}
Q_{18}={}&-82 - 414x_{8} - 529x_{8}^{2} - 1472x_{8} x_{9} - 368x_{8} x_{19} + 92x_{8} x_{20} - 576x_{9}\\
&\quad -1024x_{9}^{2} - 512x_{9} x_{19} + 128x_{9} x_{20} - 144x_{19} - 64x_{19}^{2} + 32x_{19} x_{20} + 36x_{20}\\
&\quad -4x_{20}^{2} + x_{33}
\end{split}\end{equation*}
\begin{equation*}\begin{split}
Q_{19}={}&-4159 - 23994x_{8} - 34596x_{8}^{2} - 95232x_{8} x_{9} - 23808x_{8} x_{19} + 5952x_{8} x_{20} + 744x_{8} x_{33}\\
&\quad -33024x_{9} - 65536x_{9}^{2} - 32768x_{9} x_{19} + 8192x_{9} x_{20} + 1024x_{9} x_{33} - 8256x_{19} - 4096x_{19}^{2}\\
&\quad + 2048x_{19} x_{20} + 256x_{19} x_{33} + 2064x_{20} - 256x_{20}^{2} - 64x_{20} x_{33} + 258x_{33} - 4x_{33}^{2}\\
&\quad + x_{34}
\end{split}\end{equation*}
\begin{equation*}\begin{split}
Q_{20}={}&-x_{9}^{2} - 4x_{9} x_{34} - 4x_{34}^{2} + x_{35}
\end{split}\end{equation*}
\begin{equation*}\begin{split}
Q_{21}={}&-6 - 30x_{8} - 64x_{8}^{2} - 128x_{8} x_{9} - 64x_{8} x_{19} - 32x_{9} - 64x_{9}^{2}\\
&\quad -64x_{9} x_{19} - 14x_{19} - 16x_{19}^{2} + x_{36}
\end{split}\end{equation*}
\begin{equation*}\begin{split}
Q_{22}={}&-4 + 8x_{9} - 64x_{9}^{2} - 96x_{9} x_{10} - 32x_{9} x_{23} - 32x_{9} x_{36} + 6x_{10}\\
&\quad -36x_{10}^{2} - 24x_{10} x_{23} - 24x_{10} x_{36} + 2x_{23} - 4x_{23}^{2} - 8x_{23} x_{36} + 4x_{36}\\
&\quad -4x_{36}^{2} + x_{37}
\end{split}\end{equation*}
\begin{equation*}\begin{split}
Q_{23}={}&8x_{8} - 64x_{8}^{2} - 96x_{8} x_{10} - 32x_{8} x_{24} - 32x_{8} x_{37} + 6x_{10} - 36x_{10}^{2}\\
&\quad -24x_{10} x_{24} - 24x_{10} x_{37} + 2x_{24} - 4x_{24}^{2} - 8x_{24} x_{37} + 4x_{37} - 4x_{37}^{2}\\
&\quad + x_{38}
\end{split}\end{equation*}
\begin{equation*}\begin{split}
Q_{24}={}&-x_{10}^{2} - 2x_{10} x_{11} - x_{11}^{2} + x_{39}
\end{split}\end{equation*}
\begin{equation*}\begin{split}
Q_{25}={}&-35 + 26x_{8} - 5x_{8}^{2} - 8x_{8} x_{9} - 4x_{8} x_{10} - 8x_{8} x_{23} - 8x_{8} x_{24}\\
&\quad -4x_{8} x_{39} + 24x_{9} - 4x_{9}^{2} - 8x_{9} x_{23} - 8x_{9} x_{24} + 4x_{10} - 4x_{10}^{2}\\
&\quad -8x_{10} x_{39} + 24x_{23} - 4x_{23}^{2} - 8x_{23} x_{24} + 24x_{24} - 4x_{24}^{2} + 4x_{39}\\
&\quad -4x_{39}^{2} + x_{40}
\end{split}\end{equation*}
\begin{equation*}\begin{split}
Q_{26}={}&4x_{8} - 18x_{8}^{2} - 40x_{8} x_{10} - 32x_{8} x_{24} - 40x_{8} x_{39} - 16x_{8} x_{40} + 4x_{10}\\
&\quad -24x_{10}^{2} - 32x_{10} x_{24} - 48x_{10} x_{39} - 16x_{10} x_{40} - 16x_{24}^{2} - 32x_{24} x_{39} - 16x_{24} x_{40}\\
&\quad + 4x_{39} - 24x_{39}^{2} - 16x_{39} x_{40} - 4x_{40}^{2} + x_{41}
\end{split}\end{equation*}
\begin{equation*}\begin{split}
Q_{27}={}&x_{10} - 25x_{10}^{2} - 40x_{10} x_{24} - 20x_{10} x_{27} - 16x_{24}^{2} - 16x_{24} x_{27} - 4x_{27}^{2}\\
&\quad + x_{42}
\end{split}\end{equation*}
\begin{equation*}\begin{split}
Q_{28}={}&-1 + 8x_{9} - 4x_{9}^{2} - 8x_{9} x_{23} + 8x_{23} - 4x_{23}^{2} + x_{43}
\end{split}\end{equation*}
\begin{equation*}\begin{split}
Q_{29}={}&-1 - 16x_{9}^{2} - 8x_{9} x_{10} - 32x_{9} x_{23} - 16x_{9} x_{43} - x_{10}^{2} - 8x_{10} x_{23}\\
&\quad -4x_{10} x_{43} - 16x_{23}^{2} - 16x_{23} x_{43} - 4x_{43}^{2} + x_{44}
\end{split}\end{equation*}
\begin{equation*}\begin{split}
Q_{30}={}&1 + x_{9} - x_{9}^{2} - 4x_{9} x_{10} - 4x_{9} x_{44} - 4x_{10}^{2} - 8x_{10} x_{44}\\
&\quad -4x_{44}^{2} + x_{45}
\end{split}\end{equation*}
\begin{equation*}\begin{split}
Q_{31}={}&-1 - 4x_{45} - 4x_{45}^{2} + x_{46}
\end{split}\end{equation*}
\begin{equation*}\begin{split}
Q_{32}={}&-1 - x_{10}^{2} - 8x_{10} x_{45} - 4x_{10} x_{46} - 16x_{45}^{2} - 16x_{45} x_{46} - 4x_{46}^{2}\\
&\quad + x_{47}
\end{split}\end{equation*}
\begin{equation*}\begin{split}
Q_{33}={}&2 + x_{9} - x_{9}^{2} - 4x_{9} x_{10} - 4x_{9} x_{47} - 4x_{10}^{2} - 8x_{10} x_{47}\\
&\quad -4x_{47}^{2} + x_{48}
\end{split}\end{equation*}
\begin{equation*}\begin{split}
Q_{34}={}&-4x_{42}^{2} - 8x_{42} x_{48} - 4x_{48}^{2} + x_{49}
\end{split}\end{equation*}
\begin{equation*}\begin{split}
Q_{35}={}&-9 + 7x_{8} - x_{8}^{2} - 8x_{8} x_{42} - 4x_{8} x_{49} + 24x_{42} - 16x_{42}^{2}\\
&\quad -16x_{42} x_{49} + 12x_{49} - 4x_{49}^{2} + x_{50}
\end{split}\end{equation*}
\begin{equation*}\begin{split}
Q_{36}={}&-1 - 4x_{50} - 4x_{50}^{2} + x_{51}
\end{split}\end{equation*}
\begin{equation*}\begin{split}
Q_{37}={}&-x_{10}^{2} - 8x_{10} x_{50} - 4x_{10} x_{51} - 16x_{50}^{2} - 16x_{50} x_{51} - 4x_{51}^{2} + x_{52}
\end{split}\end{equation*}
\begin{equation*}\begin{split}
Q_{38}={}&-x_{8} - x_{8}^{2} - 4x_{8} x_{10} - 4x_{8} x_{52} - 4x_{10} - 4x_{10}^{2} - 8x_{10} x_{52}\\
&\quad -4x_{52} - 4x_{52}^{2} + x_{53}
\end{split}\end{equation*}
\begin{equation*}\begin{split}
Q_{39}={}&-10 - 12x_{53} - 4x_{53}^{2} + x_{54}
\end{split}\end{equation*}
\begin{equation*}\begin{split}
Q_{40}={}&1 + 4x_{9} - 4x_{9}^{2} - 8x_{9} x_{23} - 16x_{9} x_{53} - 8x_{9} x_{54} + 4x_{23}\\
&\quad -4x_{23}^{2} - 16x_{23} x_{53} - 8x_{23} x_{54} + 8x_{53} - 16x_{53}^{2} - 16x_{53} x_{54} + 4x_{54}\\
&\quad -4x_{54}^{2} + x_{55}
\end{split}\end{equation*}
\begin{equation*}\begin{split}
Q_{41}={}&-10 + x_{9} - 25x_{9}^{2} - 40x_{9} x_{23} - 20x_{9} x_{55} - 16x_{23}^{2} - 16x_{23} x_{55}\\
&\quad -12x_{48} - 4x_{48}^{2} - 4x_{55}^{2} + x_{56}
\end{split}\end{equation*}
\begin{equation*}\begin{split}
Q_{42}={}&21 + 24x_{48} - 8x_{48}^{2} - 16x_{48} x_{56} - 2x_{56} - 4x_{56}^{2} + x_{57}
\end{split}\end{equation*}
\begin{equation*}\begin{split}
Q_{43}={}&1 + 4x_{8} - 4x_{8}^{2} - 8x_{8} x_{19} - 8x_{8} x_{59} + x_{10} - x_{10}^{2}\\
&\quad -4x_{10} x_{57} + 4x_{19} - 4x_{19}^{2} - 8x_{19} x_{59} - 4x_{57}^{2} + x_{58} + 4x_{59}\\
&\quad -4x_{59}^{2}
\end{split}\end{equation*}
\begin{equation*}\begin{split}
Q_{44}={}&-4x_{8} - 4x_{8}^{2} - 4x_{8} x_{10} - 8x_{8} x_{19} - x_{10} - x_{10}^{2} - 4x_{10} x_{19}\\
&\quad -4x_{19} - 4x_{19}^{2} + x_{59}
\end{split}\end{equation*}
\begin{equation*}\begin{split}
Q_{45}={}&4x_{8} - 24x_{8}^{2} - 48x_{8} x_{19} - 16x_{8} x_{58} - 16x_{8} x_{59} + 4x_{19} - 24x_{19}^{2}\\
&\quad -16x_{19} x_{58} - 16x_{19} x_{59} - 4x_{58}^{2} + 8x_{59} - 8x_{59}^{2} + x_{60}
\end{split}\end{equation*}
\begin{equation*}\begin{split}
Q_{46}={}&-2 - 4x_{60} - 4x_{60}^{2} + x_{61}
\end{split}\end{equation*}
\begin{equation*}\begin{split}
Q_{47}={}&-x_{2} x_{13} + x_{62}
\end{split}\end{equation*}
\begin{equation*}\begin{split}
Q_{48}={}&x_{12} - x_{13} - 2x_{22} x_{62} - x_{63}
\end{split}\end{equation*}
\begin{equation*}\begin{split}
Q_{49}={}&x_{12} - x_{13} - 2x_{26} x_{62} - x_{64}
\end{split}\end{equation*}
\begin{equation*}\begin{split}
Q_{50}={}&x_{12} - x_{13} - 2x_{28} x_{62} - 6x_{62} - x_{65}
\end{split}\end{equation*}
\begin{equation*}\begin{split}
Q_{51}={}&x_{12} - x_{13} - 2x_{30} x_{62} - x_{66}
\end{split}\end{equation*}
\begin{equation*}\begin{split}
Q_{52}={}&x_{12} - x_{13} - 2x_{32} x_{62} + 6x_{62} - x_{67}
\end{split}\end{equation*}
\begin{equation*}\begin{split}
Q_{53}={}&x_{12} - x_{13} - 4x_{34} x_{62} - 2x_{35} x_{62} - x_{68}
\end{split}\end{equation*}
\begin{equation*}\begin{split}
Q_{54}={}&x_{12} - x_{13} - 2x_{38} x_{62} - x_{69}
\end{split}\end{equation*}
\begin{equation*}\begin{split}
Q_{55}={}&x_{12} - x_{13} - 2x_{41} x_{62} - 4x_{62} - x_{70}
\end{split}\end{equation*}
\begin{equation*}\begin{split}
Q_{56}={}&-x_{8} x_{9} + x_{8} x_{10} + x_{9} + x_{9} x_{10} - x_{10} - x_{10}^{2} + x_{71}
\end{split}\end{equation*}
\begin{equation*}\begin{split}
Q_{57}={}&1 + 2x_{11} + x_{11}^{2} - x_{71}^{2} + x_{72}
\end{split}\end{equation*}
\begin{equation*}\begin{split}
Q_{58}={}&-x_{12}^{2} + 2x_{12} x_{13} + 2x_{12} x_{74} - x_{13}^{2} - 2x_{13} x_{74} - x_{72}^{2} + x_{73}\\
&\quad -x_{74}^{2}
\end{split}\end{equation*}
\begin{equation*}\begin{split}
Q_{59}={}&-4x_{60} x_{62} - 2x_{61} x_{62} + 2x_{62} + x_{74}
\end{split}\end{equation*}
\begin{equation*}\begin{split}
Q_{60}={}&-x_{9} x_{62} - x_{13} + x_{75} x_{76} + x_{77}
\end{split}\end{equation*}
\begin{equation*}\begin{split}
Q_{61}={}&-1 + 2x_{2} - 2x_{2} x_{8} - 2x_{2} x_{19} + x_{75}
\end{split}\end{equation*}
\begin{equation*}\begin{split}
Q_{62}={}&-x_{8} x_{202} - x_{12} + x_{76} + x_{202}
\end{split}\end{equation*}
\begin{equation*}\begin{split}
Q_{63}={}&-6x_{8} x_{62} - 8x_{9} x_{62} + x_{9} x_{202} + x_{12} - x_{13} - 2x_{19} x_{62} + x_{78}
\end{split}\end{equation*}
\begin{equation*}\begin{split}
Q_{64}={}&-x_{12} x_{63} + x_{13} x_{63} + x_{79}
\end{split}\end{equation*}
\begin{equation*}\begin{split}
Q_{65}={}&-x_{64} x_{79} + x_{80}
\end{split}\end{equation*}
\begin{equation*}\begin{split}
Q_{66}={}&-x_{65} x_{80} + x_{81}
\end{split}\end{equation*}
\begin{equation*}\begin{split}
Q_{67}={}&-x_{66} x_{81} + x_{82}
\end{split}\end{equation*}
\begin{equation*}\begin{split}
Q_{68}={}&-x_{67} x_{82} + x_{83}
\end{split}\end{equation*}
\begin{equation*}\begin{split}
Q_{69}={}&-x_{68} x_{83} + x_{84}
\end{split}\end{equation*}
\begin{equation*}\begin{split}
Q_{70}={}&-x_{69} x_{84} + x_{85}
\end{split}\end{equation*}
\begin{equation*}\begin{split}
Q_{71}={}&-x_{70} x_{85} + x_{86}
\end{split}\end{equation*}
\begin{equation*}\begin{split}
Q_{72}={}&-x_{77} x_{86} + x_{87}
\end{split}\end{equation*}
\begin{equation*}\begin{split}
Q_{73}={}&-x_{78} x_{87} + x_{88}
\end{split}\end{equation*}
\begin{equation*}\begin{split}
Q_{74}={}&-x_{88}^{2} + x_{89}
\end{split}\end{equation*}
\begin{equation*}\begin{split}
Q_{75}={}&-x_{4} - 2x_{4} x_{16} + 2x_{16} + x_{90}
\end{split}\end{equation*}
\begin{equation*}\begin{split}
Q_{76}={}&-x_{14} x_{90} + x_{91}
\end{split}\end{equation*}
\begin{equation*}\begin{split}
Q_{77}={}&-2x_{1}^{2} + 2x_{1} x_{12} + 2x_{1} x_{13} + 2x_{1} x_{91} + 2x_{1} x_{92} - x_{12}^{2} - 2x_{12} x_{91}\\
&\quad -x_{13}^{2} - 2x_{13} x_{92} + x_{17} x_{94} - x_{73} x_{89} - x_{91}^{2} - x_{92}^{2}
\end{split}\end{equation*}
\begin{equation*}\begin{split}
Q_{78}={}&-x_{4} x_{15} + x_{15} - x_{15} x_{90} + x_{92}
\end{split}\end{equation*}
\begin{equation*}\begin{split}
Q_{79}={}&-x_{95} x_{98} + x_{98} + x_{100}
\end{split}\end{equation*}
\begin{equation*}\begin{split}
Q_{80}={}&-x_{3} x_{121} + x_{96} x_{194} - x_{99} - x_{100} x_{121}
\end{split}\end{equation*}
\begin{equation*}\begin{split}
Q_{81}={}&x_{101} x_{194} - x_{102} x_{124} - x_{103}
\end{split}\end{equation*}
\begin{equation*}\begin{split}
Q_{82}={}&x_{104} x_{194} - x_{105} x_{124} - x_{106}
\end{split}\end{equation*}
\begin{equation*}\begin{split}
Q_{83}={}&x_{107} x_{194} - x_{108} x_{124} - x_{109}
\end{split}\end{equation*}
\begin{equation*}\begin{split}
Q_{84}={}&x_{110} x_{194} - x_{111} x_{124} - x_{112}
\end{split}\end{equation*}
\begin{equation*}\begin{split}
Q_{85}={}&x_{113} x_{195} - x_{115} x_{124} - x_{116}
\end{split}\end{equation*}
\begin{equation*}\begin{split}
Q_{86}={}&x_{117} x_{195} - x_{118} x_{124} - x_{119}
\end{split}\end{equation*}
\begin{equation*}\begin{split}
Q_{87}={}&-x_{17} x_{121} + x_{124}
\end{split}\end{equation*}
\begin{equation*}\begin{split}
Q_{88}={}&x_{120} - x_{122} x_{124} - x_{123}
\end{split}\end{equation*}
\begin{equation*}\begin{split}
Q_{89}={}&-x_{124} x_{126} + x_{125} - x_{127}
\end{split}\end{equation*}
\begin{equation*}\begin{split}
Q_{90}={}&-x_{124} x_{129} + x_{128} x_{194} - x_{130}
\end{split}\end{equation*}
\begin{equation*}\begin{split}
Q_{91}={}&x_{131} - 2x_{131} x_{133} + x_{134}
\end{split}\end{equation*}
\begin{equation*}\begin{split}
Q_{92}={}&-x_{131} x_{134} - x_{134} + x_{135}
\end{split}\end{equation*}
\begin{equation*}\begin{split}
Q_{93}={}&2x_{132} + x_{132}^{2} + x_{134}^{2} - x_{135}^{2}
\end{split}\end{equation*}
\begin{equation*}\begin{split}
Q_{94}={}&-x_{132} x_{137} - x_{137} + x_{138}
\end{split}\end{equation*}
\begin{equation*}\begin{split}
Q_{95}={}&-1 + x_{136}^{2} + 2x_{136} x_{138} + x_{137}^{2}
\end{split}\end{equation*}
\begin{equation*}\begin{split}
Q_{96}={}&-x_{132} x_{140} - x_{140} + x_{141}
\end{split}\end{equation*}
\begin{equation*}\begin{split}
Q_{97}={}&-1 + x_{139}^{2} + 2x_{139} x_{141} + x_{140}^{2}
\end{split}\end{equation*}
\begin{equation*}\begin{split}
Q_{98}={}&1 + x_{132} - x_{132} x_{143} - x_{143} + x_{144}
\end{split}\end{equation*}
\begin{equation*}\begin{split}
Q_{99}={}&x_{142}^{2} + 2x_{142} x_{144} - 2x_{143} + x_{143}^{2}
\end{split}\end{equation*}
\begin{equation*}\begin{split}
Q_{100}={}&-x_{132} x_{146} - x_{146} + x_{147}
\end{split}\end{equation*}
\begin{equation*}\begin{split}
Q_{101}={}&-1 + x_{145}^{2} + 2x_{145} x_{147} + x_{146}^{2}
\end{split}\end{equation*}
\begin{equation*}\begin{split}
Q_{102}={}&-x_{132} x_{149} - x_{149} + x_{150}
\end{split}\end{equation*}
\begin{equation*}\begin{split}
Q_{103}={}&-1 + x_{148}^{2} + 2x_{148} x_{150} + x_{149}^{2}
\end{split}\end{equation*}
\begin{equation*}\begin{split}
Q_{104}={}&-x_{132} x_{152} - x_{152} + x_{153}
\end{split}\end{equation*}
\begin{equation*}\begin{split}
Q_{105}={}&-1 + x_{151}^{2} + 2x_{151} x_{153} + x_{152}^{2}
\end{split}\end{equation*}
\begin{equation*}\begin{split}
Q_{106}={}&-x_{132} x_{155} - x_{155} + x_{156}
\end{split}\end{equation*}
\begin{equation*}\begin{split}
Q_{107}={}&-1 + x_{154}^{2} + 2x_{154} x_{156} + x_{155}^{2}
\end{split}\end{equation*}
\begin{equation*}\begin{split}
Q_{108}={}&-x_{132} x_{158} - x_{158} + x_{159}
\end{split}\end{equation*}
\begin{equation*}\begin{split}
Q_{109}={}&-1 + x_{157}^{2} + 2x_{157} x_{159} + x_{158}^{2}
\end{split}\end{equation*}
\begin{equation*}\begin{split}
Q_{110}={}&-x_{132} x_{161} - x_{161} + x_{162}
\end{split}\end{equation*}
\begin{equation*}\begin{split}
Q_{111}={}&-1 + x_{160}^{2} + 2x_{160} x_{162} + x_{161}^{2}
\end{split}\end{equation*}
\begin{equation*}\begin{split}
Q_{112}={}&-x_{132} x_{164} - x_{164} + x_{165}
\end{split}\end{equation*}
\begin{equation*}\begin{split}
Q_{113}={}&-1 + x_{163}^{2} + 2x_{163} x_{165} + x_{164}^{2}
\end{split}\end{equation*}
\begin{equation*}\begin{split}
Q_{114}={}&1 - x_{132} x_{167} - x_{167} + x_{168}
\end{split}\end{equation*}
\begin{equation*}\begin{split}
Q_{115}={}&-1 + 2x_{166} + x_{166}^{2} + 2x_{166} x_{168} + x_{167}^{2}
\end{split}\end{equation*}
\begin{equation*}\begin{split}
Q_{116}={}&-x_{132} x_{170} - x_{170} + x_{171}
\end{split}\end{equation*}
\begin{equation*}\begin{split}
Q_{117}={}&-1 + x_{169}^{2} + 2x_{169} x_{171} + x_{170}^{2}
\end{split}\end{equation*}
\begin{equation*}\begin{split}
Q_{118}={}&-x_{132} x_{173} - x_{173} + x_{174}
\end{split}\end{equation*}
\begin{equation*}\begin{split}
Q_{119}={}&-1 + x_{172}^{2} + 2x_{172} x_{174} + x_{173}^{2}
\end{split}\end{equation*}
\begin{equation*}\begin{split}
Q_{120}={}&-x_{132} x_{176} - x_{176} + x_{177}
\end{split}\end{equation*}
\begin{equation*}\begin{split}
Q_{121}={}&-1 + x_{175}^{2} + 2x_{175} x_{177} + x_{176}^{2}
\end{split}\end{equation*}
\begin{equation*}\begin{split}
Q_{122}={}&-5 - 2x_{0} - 4x_{0}^{2} - 2x_{0} x_{2} - x_{1}^{2} - 2x_{2}^{2} - x_{3}^{2}\\
&\quad -x_{4}^{2} + 2x_{8} - x_{8}^{2} - x_{9}^{2} - x_{10}^{2} - x_{11}^{2} - x_{12}^{2}\\
&\quad -x_{13}^{2} - x_{14}^{2} - x_{15}^{2} - 4x_{16}^{2} - x_{17}^{2} - x_{93}^{2} + 2x_{95}\\
&\quad -2x_{95}^{2} - x_{96}^{2} - x_{97}^{2} - x_{101}^{2} - x_{104}^{2} - x_{107}^{2} - x_{110}^{2}\\
&\quad -x_{113}^{2} - x_{114}^{2} - x_{117}^{2} - x_{120}^{2} - x_{121}^{2} - x_{125}^{2} - x_{128}^{2}\\
&\quad + x_{131} - x_{178}^{2}
\end{split}\end{equation*}
\begin{equation*}\begin{split}
Q_{123}={}&2x_{3} - 4x_{3} x_{8} - 4x_{3} x_{9} - 4x_{3} x_{10} - 4x_{3} x_{11} - 4x_{3} x_{12} - 4x_{3} x_{13}\\
&\quad -4x_{3} x_{14} - 4x_{3} x_{15} - 8x_{3} x_{16} + x_{178}
\end{split}\end{equation*}
\begin{equation*}\begin{split}
Q_{124}={}&-128 + x_{0} - 4x_{179}^{2}
\end{split}\end{equation*}
\begin{equation*}\begin{split}
Q_{125}={}&-x_{132} x_{181} - x_{181} + x_{190}
\end{split}\end{equation*}
\begin{equation*}\begin{split}
Q_{126}={}&-1 + x_{180}^{2} + x_{181}^{2} - x_{190}^{2}
\end{split}\end{equation*}
\begin{equation*}\begin{split}
Q_{127}={}&-x_{183} - x_{183} x_{184} + x_{191}
\end{split}\end{equation*}
\begin{equation*}\begin{split}
Q_{128}={}&-1 + x_{182}^{2} + x_{183}^{2} - x_{191}^{2}
\end{split}\end{equation*}
\begin{equation*}\begin{split}
Q_{129}={}&-x_{173}^{2} + x_{192}
\end{split}\end{equation*}
\begin{equation*}\begin{split}
Q_{130}={}&x_{181} - 2x_{186} x_{192} - x_{192}
\end{split}\end{equation*}
\begin{equation*}\begin{split}
Q_{131}={}&-4x_{181} - 8x_{181} x_{185} + x_{184}
\end{split}\end{equation*}
\begin{equation*}\begin{split}
Q_{132}={}&-x_{132} - x_{180} x_{187} + x_{184}
\end{split}\end{equation*}
\begin{equation*}\begin{split}
Q_{133}={}&-x_{172} - x_{174} - x_{180} x_{188} + x_{182}
\end{split}\end{equation*}
\begin{equation*}\begin{split}
Q_{134}={}&-2x_{16} - 4x_{173} x_{189} + x_{183}
\end{split}\end{equation*}
\begin{equation*}\begin{split}
Q_{135}={}&26 + 2x_{0}^{2} + 2x_{0} x_{2} - 2x_{0} x_{137} - 2x_{0} x_{140} + x_{1}^{2} - 2x_{1} x_{167}\\
&\quad + x_{2}^{2} - 2x_{2} x_{140} + x_{4}^{2} - 2x_{4} x_{158} + x_{8}^{2} - 2x_{8} x_{143} + x_{9}^{2}\\
&\quad -2x_{9} x_{146} + x_{10}^{2} - 2x_{10} x_{149} + x_{11}^{2} - 2x_{11} x_{152} + x_{12}^{2} - 2x_{12} x_{155}\\
&\quad + x_{13}^{2} - 2x_{13} x_{161} + x_{14}^{2} - 2x_{14} x_{164} + x_{15}^{2} - 2x_{15} x_{170} + 4x_{16}^{2}\\
&\quad -4x_{16} x_{173} + x_{93}^{2} - 2x_{93} x_{176} - x_{132} x_{193} - 2x_{136} - 2x_{138} + x_{138}^{2}\\
&\quad -2x_{139} - 2x_{141} + x_{141}^{2} - 2x_{142} + 2x_{143} - 2x_{144} + x_{144}^{2}\\
&\quad -2x_{145} - 2x_{147} + x_{147}^{2} - 2x_{148} - 2x_{150} + x_{150}^{2} - 2x_{151}\\
&\quad -2x_{153} + x_{153}^{2} - 2x_{154} - 2x_{156} + x_{156}^{2} - 2x_{157} - 2x_{159}\\
&\quad + x_{159}^{2} - 2x_{160} - 2x_{162} + x_{162}^{2} - 2x_{163} - 2x_{165} + x_{165}^{2}\\
&\quad -2x_{166} + x_{168}^{2} - 2x_{169} - 2x_{171} + x_{171}^{2} - 2x_{172} - 2x_{174}\\
&\quad + x_{174}^{2} - 2x_{175} - 2x_{177} + x_{177}^{2}
\end{split}\end{equation*}
\begin{equation*}\begin{split}
Q_{136}={}&-x_{97} x_{194} + x_{121}
\end{split}\end{equation*}
\begin{equation*}\begin{split}
Q_{137}={}&-x_{114} x_{195} + x_{121}
\end{split}\end{equation*}
\begin{equation*}\begin{split}
Q_{138}={}&-2 - 4x_{16}^{2} - x_{137}^{2} - x_{140}^{2} + 2x_{143} - x_{143}^{2} - x_{146}^{2}\\
&\quad -x_{149}^{2} - x_{152}^{2} - x_{155}^{2} - x_{158}^{2} - x_{161}^{2} - x_{164}^{2} - x_{167}^{2}\\
&\quad -x_{170}^{2} + x_{173}^{2} - x_{176}^{2} - x_{196}^{2} - x_{197}^{2} - x_{198} - x_{198}^{2}
\end{split}\end{equation*}
\begin{equation*}\begin{split}
Q_{139}={}&-1 - 4x_{99}^{2} - 4x_{103}^{2} - 4x_{106}^{2} - 4x_{109}^{2} - 4x_{112}^{2} - 4x_{116}^{2}\\
&\quad -4x_{119}^{2} + x_{121}^{2} - 4x_{123}^{2} - 4x_{127}^{2} - 4x_{130}^{2} - x_{199}^{2} - x_{200}^{2}\\
&\quad -x_{201} - x_{201}^{2}
\end{split}\end{equation*}
\begin{equation*}\begin{split}
Q_{140}={}&-x_{2} x_{12} + x_{202}
\end{split}\end{equation*}
\begin{equation*}\begin{split}
Q_{141}={}&-1 - x_{2} - 2x_{2} x_{203} + x_{4}
\end{split}\end{equation*}
\begin{equation*}\begin{split}
Q_{142}={}&-x_{0}^{2} + x_{204}
\end{split}\end{equation*}
\begin{equation*}\begin{split}
Q_{143}={}&x_{93} - x_{204}^{2}
\end{split}\end{equation*}
\endgroup



\section{Companion quartic and machine-readable identities}

$P_{228}$ is the sum of the 175 residual squares in
\texttt{equations/P228\_compact/equation.json}; every collected coefficient
is also supplied in that directory. Its readable residual list and
LaTeX include are complete. It is obtained from the compact package-166
parent by exactly the same two-row change and private variable
identification proved above. The balanced and compact parents use
different integral projections; one must not delete a coordinate from
the balanced list and assume that the compact equation follows without
its separate map.

The source-coordinate and final-coordinate certificates define every
map, including all restored eliminated variables. The sorted sparse
integer lists are authoritative mathematical data; PDF metadata and
font bytes may differ across TeX installations without changing any
equation. The final checksum manifest covers the supplied files except
itself. Historical source notices are preserved, and the new work uses
the repository's MIT No Attribution terms.

\begin{thebibliography}{9}
\raggedright
\bibitem{collection} Vladimir Reshetnikov, \emph{ProveIt: Con(ZFC) in
Diophantine form}, source snapshot \pin.
\href{https://github.com/VladimirReshetnikov/ProveIt/tree/02a78a08603cbae248c19c23a75907c6c6bd6443/Computability/DiophantineConsistency/con-zfc-diophantine-forms}{Pinned collection and provenance}.
\bibitem{source166} ProveIt, package 166, \emph{Fixed-base modular
Diophantine consistency}, 25 October 2026. In particular
\texttt{TWO\_SQUARE\_CENTER.md}, \texttt{FRAME\_AND\_AFFINE\_PROOF.md},
and \texttt{MODULAR\_PROOF.md}, Section 10.
\href{https://github.com/VladimirReshetnikov/ProveIt/tree/02a78a08603cbae248c19c23a75907c6c6bd6443/Computability/DiophantineConsistency/con-zfc-diophantine-forms/166-fixed-base-modular}{Pinned package 166}.
\bibitem{source168} ProveIt, package 168, \emph{Reciprocal pooling and
integral projections for Con(ZFC)}, 25 October 2026.
\href{https://github.com/VladimirReshetnikov/ProveIt/tree/02a78a08603cbae248c19c23a75907c6c6bd6443/Computability/DiophantineConsistency/con-zfc-diophantine-forms/168-reciprocal-pooling}{Pinned package 168}.
\bibitem{source173} ProveIt, package 173, \emph{Reusing the encoded row
index as the Pell anchor}, 25 October 2026. Source chart
\texttt{source/anchor221/equation.json}; proofs
\texttt{C\_ANCHOR\_THEOREM.md},
\texttt{LABEL\_SEPARATION\_WITHOUT\_DERIVED\_RANGE.md}, and
\texttt{ORDINARY\_PROOF\_DEPENDENCY\_BOUNDARY.md}.
\href{https://github.com/VladimirReshetnikov/ProveIt/tree/02a78a08603cbae248c19c23a75907c6c6bd6443/Computability/DiophantineConsistency/con-zfc-diophantine-forms/173-crt-index-anchor-215}{Pinned package 173}.
\bibitem{endpoint} ProveIt, \texttt{KEndpoint.lean}, theorem
\texttt{c222\_iff\_zf\_refutation}.
\href{https://github.com/VladimirReshetnikov/ProveIt/blob/02a78a08603cbae248c19c23a75907c6c6bd6443/Computability/BusyBeaver/Lean/BusyBeaver/ZFC/KEndpoint.lean}{Pinned endpoint source}.
\bibitem{choice} ProveIt, \texttt{ConstructibleChoice.lean}, theorem
\texttt{zf\_refutation\_iff\_zfc\_refutation}.
\href{https://github.com/VladimirReshetnikov/ProveIt/blob/02a78a08603cbae248c19c23a75907c6c6bd6443/Computability/BusyBeaver/Lean/BusyBeaver/ZFC/ConstructibleChoice.lean}{Pinned Choice transport source}.
\bibitem{sun} Zhi-Wei Sun, \emph{Ten unknowns for Hilbert's tenth problem
over the integers}, arXiv:2609.01594, 2026.
\url{https://arxiv.org/abs/2609.01594}.
\end{thebibliography}
\end{document}
